\documentclass[../main.tex]{subfiles}
\begin{document}
%==========================================TIÊU ĐỀ TẬP========================================
\begin{table}[h]
  \begin{adjustwidth}{-1cm}{}
    \begin{tabular}{>{\centering\arraybackslash}p{6cm}|>{\raggedright\arraybackslash}p{14cm}}
      \multirow{5}{6cm}
      % Cột bên trái: ÉPISODE và số tập
      {\\[-40pt]
        \flaregothic\fontsize{20pt}{20pt}\selectfont \centering
        \color{eptype}HISTOIRE\color{black}\\
        \gangofthree\fontsize{72pt}{72pt}\selectfont
        \color{epnum}9
        \\[18pt]
        % \flaregothic\fontsize{14pt}{14pt}\selectfont
        % \color{epattr}BIENVENUE, \uline{\color{Banpire} Mel} !
        % \flaregothic\fontsize{18pt}{18pt}\selectfont
        % \color{epattr}ÉPISODE PERDU
        \color{black}
      }
       &
      % Dòng thứ nhất: tên tập tiếng Nhật
      {\fotskip\fontsize{18pt}{18pt}\selectfont \ruby{正}{しょう}\ruby{月}{がつ}は\ruby{希}{き}\ruby{望}{ぼう}の\ruby{象}{しょう}\ruby{徴}{ちょう}}
      \\[6pt]
      % Dòng thứ hai: tên tập tiếng Anh
       & {\cafeteria\fontsize{18pt}{18pt}\selectfont Tet Means Hope}                       \\[4pt]
      % Dòng thứ ba: tên tập tiếng Việt
       & {\vnmsans\fontsize{18pt}{18pt}\selectfont Tết nghĩa là hy vọng}                   \\[8pt]
      % Dòng thứ tư: Nhân vật xuất hiện
       & {\brian{\fontsize{14pt}{14pt}\selectfont Track used: Lunar Zoo Year - Zombotany}}
    \end{tabular}
  \end{adjustwidth}
\end{table}
%===========================================NỘI DUNG==========================================
\color{black}

\pov{Hikawa Kuon}

Một buổi sáng rằm tháng Giêng, tôi cùng A-san ngồi lại văn phòng để chỉnh sửa nốt những
đoạn phim cuối cùng của bộ phim đặc biệt mà chúng tôi đã quay trong dịp Tết Nguyên
Đán ở Etsunan. Sau nhiều ngày liên tục làm việc, chúng tôi đã gần hoàn thành, và chỉ còn
những bước cuối cùng để mọi thứ trở nên hoàn hảo.

Vài tiếng sau, tôi gọi tất cả mọi người trong nhóm lại để xem thử đoạn phim mà cả đội đã
góp công quay và chỉnh sửa.

Tất nhiên, vì đây là bản đã sửa lần đầu, nên sẽ không tránh khỏi những sai sót. Dù vậy,
tôi vẫn quyết định cho mọi người xem thử những hoạt động trong năm ngày Tết đó được
đúc kết lại trong một bộ phim dài ba tiếng này sẽ ra sao.

``Anh muốn xem cảm nghĩ của mọi người về bộ phim này như thế nào,'' tôi nói, nhấn nút
bật máy chiếu.
Đoạn phim bắt đầu với dòng chữ lớn hiện trên màn hình:

\begin{center}
  \uline{Chán Shougatsu rồi ư?\\Vậy sao không thử ngày Tết Nguyên Đán ở Etsunan nhỉ?}
\end{center}

Sora-chan bật cười nhẹ, nhưng rồi nhanh chóng lên tiếng: ``Nếu chỉ xem tới đây thì em
nghĩ anh đang hơi làm quá rồi đấy\dots''

Subaru-chan chen vào, khuôn mặt lộ rõ vẻ khó hiểu: ``Em không nghĩ là sẽ có ai chán
ngày Tết ở nước mình đâu!''

\dots Có lẽ tôi sẽ phải thay đổi phần mở đầu thật. Ngẫm nghĩ lại, tôi có cảm giác đang hơi
xúc phạm ngày Tết cổ truyền của Nhật Bản.

Mọi người tiếp tục dõi theo. Màn hình chuyển sang dòng chữ mới, đầy màu sắc và ấn
tượng:

\begin{center}
  \uline{\hll\ the Movie\\Đại náo ngày Tết ở Etsunan!}
\end{center}

Đột nhiên, một cây giáo mác từ đâu đó xuất hiện, lao thẳng vào dòng chữ, khiến toàn bộ
tiêu đề nổ tung thành hàng loạt mảnh vụn.

Giọng A-san từ đoạn phim vang lên, rõ ràng và đầy vẻ bất bình: ``\textbf{Này! Bọn tôi mất công
  lắm mới dựng được cái tiêu đề ấy đấy!}''

Ngay sau đó, hình ảnh Noel trong phim xuất hiện, đang ôm đầu và xin lỗi rối rít: ``Em xin
lỗi chị mà~!''

Bản thân nàng Hiệp sĩ khi xem phân cảnh của mình cũng chỉ biết ôm mặt than thở: ``Sao
mình cảm thấy đang bị bêu rếu nhỉ\dots''

Tôi tạm dừng đoạn phim lại, nhìn mọi người một lượt. Bầu không khí trong phòng tràn
ngập sự hài hước pha lẫn chút ngượng ngùng của vài thành viên.

Sau một đoạn mở đầu đầy hài hước, tôi cùng A-san tiếp tục dẫn dắt mọi người vào phần
chính của bộ phim. Màn hình chiếu cảnh văn phòng, nơi các thần tượng đang ngồi lại với
nhau, bắt đầu tranh cãi về nội dung bộ phim họ sẽ làm.

\il

Matsuri, vốn không bỏ qua bất kỳ cơ hội nào, chen ngang với nụ cười đắc thắng: ``Vậy là
ý tưởng của tôi thắng rồi chứ gì? Kiểu, mọi người sẽ coi tôi như là công-''

Nhưng, cô chưa kịp nói hết câu thì Fubuki thản nhiên chen vào, giọng đầy ngạc nhiên giả
tạo: ``Ủa, bà ở đâu mọc ra đây vậy, Matsuri-chan?'' Nàng Cổ động sững lại, mắt mở to
như không tin vào tai mình: ``Nãy giờ tôi ở đây mà! Từ nãy đến giờ ta nói chuyện mà
không để ý à!?''

Bên cạnh đó, nàng Cáo trắng vẫn không ngừng bám lấy ý tưởng cũ, quay sang người bạn
của cô với ánh mắt đầy quyết tâm: ``Tôi nói rồi, chúng ta cứ làm một tập về đồ tắm đi!
Khán giả sẽ thích mà!''

Nàng Sói đen, rõ ràng đã cạn kiệt sức lực, lắc đầu và than thở: ``Và tôi đã bảo chúng ta đã
làm nhiều về nó rồi! Với cả, ai lại muốn mặc đồ bơi vào mùa đông cơ chứ!?''

Ayame, vốn dĩ lúc nào cũng góp ý cực đoan, lại chen vào: ``Vậy thì chiếm đài truyền hình
cũng được ha!'' Choco vẫn giữ vẻ mặt nghiêm túc, đưa ra một ý kiến khác: ``Cô nghĩ là
chúng ta nên theo đuổi ý tưởng về \hll\ 100\%.''

Mio gần như hét lên: ``Ta không thể làm về chúng được! Với lại, từ từ đã, cô lại đổi tiêu
đề phim nữa sao?''

Khi cuộc cãi vã bắt đầu lan sang cả những chuyện không liên quan, không khí trong phòng
ngày càng trở nên hỗn loạn.

\il

Nhìn lại cảnh mà mọi người tranh cãi nhau, tất cả các cô gái cũng chỉ cười trừ.

Fubuki-chan nhìn chính mình trong đó, mơ mộng: ``Tôi vẫn ước rằng chúng ta có thể làm
một tập về áo tắm nữa\dots'' Mio-chan nhăn mày: ``Đã bảo là chúng ta đã làm về nó rồi mà.
Đừng có mơ tưởng nữa.''

``Với cả ấy,'' đến lượt tôi lên tiếng, ``thời tiết kiểu này mà mặc áo tắm thì không phù hợp
sao? Cứ nhìn gương Mat-chan với Shion-chan hôm thi trượt tuyết là đủ hiểu rồi đấy\footnote{Xem lại {\sen Match 1} trong {\flaregothic SPÉCIAL IV}, quyển số Ba}.''

Phù thủy-chan nghe tôi nói vậy, cũng phải gật đầu đồng ý: ``Em vẫn còn nhớ cảm giác
mặc đồ tắm trên đỉnh núi hôm đó\dots''

Đoạn phim ghi lại hình ảnh từng người đưa ra ý kiến của mình. Bối cảnh lúc đó khá hỗn
loạn, nhưng tràn đầy tiếng cười. Để rồi mọi người đều bối rối với câu hỏi chủ đạo ngay
sau khi cuộc tranh luận kết thúc:

\begin{center}
  \uline{Tết là gì?}
\end{center}

Mel-chan vang lên, mẩm nhớ lại hồi mà mọi người biết tin tôi nghỉ việc để ăn Tết: ``Chị
nhớ hồi đó mọi người lo lắng lắm, tưởng Kuon-senpai bị nghỉ việc rồi mà.''

Tất cả chúng tôi đều bật cười. Rõ ràng là tôi nhớ tôi đã để lại giấy nhắn rồi mà nhỉ?

Pekora chen vào, giọng điệu như đang cố nhấn mạnh một chuyện nghiêm trọng: ``Mọi
người sợ đến xanh mặt luôn, peko.''

\il

[\dots] Roboco phá tan bầu không khí tĩnh lặng bằng việc thử đoán về định nghĩa của nó: ``Chị nghĩ chắc là một vocaloid nào đó chăng?''

``Kiểu. Kasane Teto\footnote{Trong tiếng Nhật, ``Tết'' (テト) đọc là ``teto''. Kasane Teto mà họ nhắc tới cũng là người hát bài \textit{Ochame Kinou} từ năm 2010, vốn được \hll\ cover lại vào 4/2020.} à?'' Shion hỏi lại, nheo mắt nghi ngờ, ``Em thấy hai cái này chẳng ăn khớp nhau lắm ấy.''

``Chị nghĩ là vậy\dots?'' nàng Người máy trả lời, vẻ mặt không mấy chắc chắn.

Nhắc đến vocaloid tóc đỏ hai bên xoắn ốc, Aqua và Matsuri bỗng đứng bật dậy đồng thanh hát bài \textit{Ochame Kinou}, bài hát mà cả hai vừa thu âm để mừng kỷ niệm 10 năm.

\begin{dialogue}
  \speak{Aqua} いつでも I love you, \ruby{君}{きみ}に Take kiss me{\fotskip \~{}} \\ (Lúc nào tớ cũng love you [yêu cậu], nên cậu hãy mau take [nhận lấy] rồi kiss me [hôn tớ đi]\~{})
  \speak{Matsuri} \ruby{忘}{わす}れられないから　\ruby{僕}{ぼく}の\ruby{大}{だい}\ruby{事}{じ}な\ruby{メ}{め}\ruby{モ}{も}\ruby{リ}{り}ー{\fotskip \~{}} \\ (Tớ không thể quên được những ký ức quý giá này đâu\~{})
\end{dialogue}

Mio, người vừa ngồi xuống để viết vài ghi chú, nghe thấy liền đổ mồ hôi hột: ``Tôi không nghĩ chỉ vì một người mà anh ấy sẵn sàng nghỉ gần một tháng đâu. Nó hẳn phải là một thứ gì đó to lớn hơn thế\dots''

\il

Đến phân cảnh mà Mat-chan và Akutan hát vang bài của hát của nữ vocaloid tóc đỏ đó,
tất cả chúng tôi không nhịn được, bật cười ầm lên.

Matsuri-chan và Aqua-chan đỏ mặt, đứng phắt dậy huơ huơ tay: ``M-Mọi người đừng có
cười chứ!''

Tôi vừa cười vừa nhìn cô nàng Lễ hội đang đỏ mặt kia: ``Thành thật mà nói, anh cũng
không ngờ em nghĩ đến cái tên đó đấy\dots''

Mio-chan cố gắng kiềm lại tiếng cười, nhưng rồi vẫn thẳng thắn thừa nhận: ``Cái này cũng
dễ hiểu thôi. Đến cả chúng ta còn chưa mường tượng được nó là gì mà.''

Tôi tạm dừng đoạn phim tại đây để mọi người lấy lại bình tĩnh. Trong lúc đó, A-san
liếc nhìn tôi, mỉm cười: ``Công nhận, ngay từ lúc đầu, nhóm mình đã không khác gì một
chương trình tấu hài cả.''

Tôi chỉ biết gật đầu đồng tình, còn Matsuri-chan thì ngồi im lặng, hai tay khoanh trước
ngực như đang muốn giấu đi sự ngại ngùng.

Tôi bấm nút tiếp tục chiếu phim. Sau những tranh cãi và hiểu lầm ban đầu, bộ phim dẫn
dắt mọi người đến cảnh các cô gái trong nhóm bắt đầu tìm hiểu về ngày Tết Nguyên Đán.
Qua những đoạn ghi lại các hoạt động chuẩn bị, cảnh quay chuyển sang nhà tôi, nơi cả
nhóm tụ tập lần đầu tiên.

Không gian trong phim là một ngôi nhà Etsunan nhỏ gọn nhưng ấm cúng. Tuy nhiên, cảm
giác đầu tiên của mọi người khi đến đây, như được thể hiện qua lời thoại, lại có chút\dots\ khác biệt.

Shion-chan trong phim cười khúc khích, giọng điệu nửa trêu chọc nửa chân thật: ``Nhà
anh trông có vẻ nghèo với chật chội nhỉ.''

Bên ngoài màn hình, cô nàng lại tiếp tục bồi thêm: ``Không biết bây giờ có khác hơn
không, nhưng lúc đó đúng là hơi bí thật đấy!

Tôi chỉ thở dài, cố giữ bình tĩnh mà đáp: ``\dots không phải nhắc lại đâu\dots\ Nhà anh khi
ấy còn chưa sửa sang lại xong mà.''

Subaru-chan bật cười lớn, vỗ vai tôi: ``Không sao đâu mà, Kuon-senpai! Như em ở trong
phim đã nói rồi đấy\dots''

Em ấy, cùng với bản thân bên trong màn ảnh đều cùng nhau cất lên: ``Nhà bé thì càng ấm
cúng hơn chứ, phải không?''

A-san, với vẻ mặt điềm tĩnh, lại nhìn vào màn hình và góp ý thêm: ``Tôi nghĩ không gian
như vậy tạo cảm giác rất gần gũi, đặc biệt khi so với những ngôi nhà hiện đại. Hơn nữa,
cách sắp xếp của cậu hồi đó khá chu đáo mà.''

Mọi người gật gù đồng tình với nhận xét của A-san, còn tôi chỉ biết cười trừ. Đúng là lúc
ấy, tôi đã lo lắng đến phát mệt vì phải chuẩn bị chỗ ở cho cả nhóm.

Trên màn hình, cảnh quay chuyển tiếp đến phần các cô gái bắt đầu làm quen với không
gian nhà tôi, từ việc cẩn thận di chuyển trong phòng khách nhỏ đến việc cùng nhau dọn dẹp và chuẩn bị. Tôi nghe thấy tiếng cười râm ran của các đồng nghiệp xung quanh khi
nhớ lại những khoảnh khắc hỗn loạn ấy.

Sau phần giới thiệu căn nhà cùng một chút gia thế trong gia đình Hikawa chúng tôi, bộ
phim chuyển đến cảnh mọi người được cả gia đình tôi phân công lau dọn nhà cửa để chuẩn
bị cho Tết.

Trên màn hình, từng thành viên được giao các nhiệm vụ khác nhau, từ quét nhà, mua đồ,
đến việc trang trí lại phòng ốc.

Tôi nhớ rất rõ, việc này ban đầu tưởng chừng đơn giản nhưng hóa ra lại thành một buổi
``trải nghiệm hỗn loạn.'' Và tất nhiên, tất cả mọi người đều không quên ghi lại những
khoảnh khắc đáng nhớ nhất.

\il

[\dots] Ayame bỗng thét lên khi nhặt được một vật thể đáng sợ: ``Ôi giời ơi, còn có cả gián (đã
che) này! Khiếp quá!''

Tiếng hét của nàng Quỷ sừng khiến cả Choco và Mio đứng hình trong vài giây. Khi họ
nhận ra Ayame thực sự đang cầm một con gián, cả hai đồng loạt hét lên đầy hoảng hốt:
``Sao bà/em lại cầm nó lên!?''

Nàng Quỷ sừng bình thản đáp, mặt không chút biểu cảm: ``Còn có cả một tổ gián dưới
này nữa.'' Nghe vậy, cả hai cô gái càng khiếp sợ hơn, la hét đầy hoảng loạn.

Mio ôm mặt thét: ``Éc! Giết nó đi! Giết nó ngay đi!'', còn Choco thì lùi lại, giọng run rẩy:
``Bỏ nó ra đi! Khiếp quá!''

\il

Ojou-san nhìn lại bản thân trong thước phim, mặt không giấu nổi sự hứng thứ: ``À, cái
này là lúc tôi bốc gián lên này.''
Ngay lập tức, Choco-sensei hét lên hoảng sợ, giọng run rẩy: ``Đừng nhắc lại nữa\dots\ cô sợ
lắm\dots''

Mio-chan thậm chí còn hoang mang: ``Ai đã quay lại toàn bộ cảnh này vậy!?''

Nhìn cảnh đó, bên ngoài, mọi người cười nghiêng ngả (ngoại trừ bộ ba đã dọn nhà bếp
ngày đó, khi họ ôm chầm vào nhau run rẩy), trong khi tôi thì chỉ biết thở dài: ``Đúng là
dọn nhà mà cứ như đi vào bãi chiến trường vậy\dots''

Chuyển qua một khung hình khác, Fubuki-chan trong đoạn phim đang cầm một chiếc
chổi, đôi mắt ngơ ngác nhìn xuống mảnh vỡ trên sàn nhà. Bên cạnh, Aqua-chan cuống
quýt thu dọn nhằm phi tang chứng cứ.

\il

\textbf{*XOẢNG!*}

[\dots] Tiếng kính vỡ bất ngờ vang lên. Nàng Hầu gái vừa lau dọn những chiếc cốc trên kệ thì
bất cẩn trượt tay làm rơi một chiếc xuống đất. Cô hoảng hốt quỳ xuống thu nhặt những mảnh vỡ: ``Thôi chết, mình làm vỡ cốc rồi!''

Tiếng động cũng làm cho hai người kia giật mình, vội chạy đến hỏi han. Nàng Pháp sư
lên tiếng: ``Chị có sao không!?''

Aqua đáp, vẫn vội vàng gom lại những vết nứt vỡ: ``Chị ổn\dots\ Nhưng cái cốc này bị vỡ
có sao không nhỉ?''

Fubuki đang lau một góc khác bèn liếc mắt nhìn qua rồi nói với giọng tỉnh bơ: ``Không
sao đâu~ Dù gì nó có phải cốc nhà anh ấy đâu.''

\il

Fubuki-chan ở bên ngoài nói, giọng điệu có chút tự hào: ``Cái này là lúc Akutan làm vỡ
cốc này.''

``Kéo sau đó là những thứ đổ vỡ khác nhỉ\dots'' Ru-chan vừa ngồi xem bản thân trên màn
ảnh, vừa nhớ lại những gì mà em ấy đã trải qua.

Ngay lập tức, Akutan lên tiếng phản đối: ``Đ-Đừng có nhắc lại chứ!''

Bên ngoài, A-san cười nhẹ, bình luận: ``Chị thấy cảnh này Minato-san còn không biết xử
lý thế nào mà cứ đứng đó run mãi.''

Nàng Hầu gái càng đỏ mặt hơn, thốt lên: ``Mồ! Mọi người đừng nhắc lại nữa đi!!''

Hình ảnh tiếp theo là Okayu và Korone đang trộn sơn, chuẩn bị làm mới lại căn phòng
tầng ba.

\il

[\dots] Họ biến căn phòng thành một đấu trường thực thụ. Những vật dụng như bàn kính, rèm
cửa, và thậm chí cả nhà vệ sinh được tận dụng làm nơi ẩn náu lẫn chướng ngại vật. Không
khí trở nên sôi động hơn bao giờ hết khi từng cú quất chổi và những loạt đạn sơn bắt đầu
nhuộm bức tường đầy sắc màu.

Korone nhìn vũ khí của Okayu, tỏ vẻ bất bình: ``Ể~? Okayu không công bằng! Bà có
súng kìa!''

Okayu nở nụ cười, nhún vai đáp lại: ``Bà cũng có cả tá lựu đạn sơn đấy thôi? Với cả, cờ
tay ai người ấy phất mà?''

Nàng Khuyển liếc xuống thắt lưng của mình, nơi treo vài quả lựu đạn sơn màu cam rực
rỡ, bật cười: ``À, đúng rồi nhỉ\dots''

Nàng Miêu giơ súng lên cao, hô hào: ``Đây chỉ là một trận giao hữu thôi. Cứ thoải mái đi
nha~''

Cả căn phòng nhanh chóng trở thành chiến trường đầy ắp tiếng cười. Mục tiêu chính của
họ rõ ràng không phải là sơn tường, mà là tận hưởng niềm vui từ trò đùa tinh nghịch này.

\il

Tôi chỉ tay về phía hai người họ bên trong phim và cười nói: ``Vậy ra bức tường ở tầng
ba được sơn như vậy à?''

Okayu-chan đáp lại bằng giọng trầm, điềm tĩnh: ``Pín pon~''

Koro-chan cười tươi, vẫy vẫy chiếc đuôi: ``Bọn em đã hết sức rồi đấy, Kuon-senpai!''

Và còn rất nhiều cảnh hỗn loạn khác trong lúc chúng tôi dọn nhà. Đúng là chỉ một việc
đơn giản như thế thôi mà cũng có thể thành nguồn cảm hứng cho cả một bộ phim hài.

Sau đó, màn hình chuyển sang một cảnh đặc biệt: phân đoạn ``trừng phạt'' bốn thành viên
đến từ nhóm GAMERS. Dòng tiêu đề trên màn hình hiển thị bằng chữ đỏ rực:

\begin{center}
  \uline{Hình phạt cho hành động đáng ngờ: Ngủ cùng Kuon!}
\end{center}

\il

[\dots] ``C-Có chuyện gì vậy!?'' nàng Cáo trắng vừa hất nước khỏi mặt, vừa lắp bắp hỏi, ngơ
ngác nhìn mọi người. Nàng Sói đen thì nhìn quanh, nhưng chỉ kịp nhận ra tay chân mình
đang bị trói chặt từ phía sau. Cô hoảng hốt thốt lên: ``S-Sao tay chân của tôi lại bị trói thế
này!?''

Đáp lại sự hoang mang của bốn người, giọng nói lạnh lùng của Shion vang lên từ phía xa:
``Bốn người tỉnh chưa?''

Bốn cô gái quay đầu về phía âm thanh, thấy Shion đang đứng chống hông với vẻ mặt
không mấy hài lòng. Xung quanh, tất cả các thành viên khác của \hll\ (trừ A-chan, vì
cô chưa có mặt) đang đứng đó, mỗi người một vẻ: người buồn rầu, người bất mãn, đặc
biệt là ánh mắt bừng bừng sát khí của Rushia.

  [\dots] Mio nhíu mày, không giấu được vẻ hoang mang: ``Lại còn có hình phạt nhẹ với nặng nữa
à? Rồi trước sau là sao? Hành quyết kiểu gì kỳ lạ vậy?'' còn Okayu thì ngáp nhẹ, hỏi với
giọng tỉnh rụi: ``Vậy hình phạt đó là gì thế?''

Đáp lại câu hỏi của hai người, tất cả những thành viên còn lại đồng thanh hét lên: ``\textbf{Xịt
  kem!}''

Chưa kịp hiểu chuyện gì, hai nàng đã bị cả đám cầm bình xịt kem trắng mà Shion chuẩn
bị trước phun thẳng vào người.

Nàng Sói đen hét toáng lên: ``Á!! Sao nhiều quá vậy!?'' Nàng Mèo thì cười khúc khích,
bất chấp tình cảnh của mình: ``Ui, cái này làm em thấy hơi buồn (nhột) nha~''

Korone, người đang bị cắm cọc ngay bên cạnh, trố mắt nhìn cô: ``Ơ\dots\ Ogayu cười vì bị
kem chọc lét à?'' Fubuki thì cười với vẻ thích thú: ``Trông có vẻ vui à nha~''

[\dots] Toàn bộ gần hai mươi người còn lại lao đến, bắt đầu cúi người cù lét khắp toàn thân của
cả bốn người.

Nàng Cáo trắng gào lên: ``\textbf{D-Dừng lại đi- Ahahaha!}'' Nàng Miêu tím thì vừa cười vừa
lẩm bẩm: ``Nữa đi, nữa đi~''

Nàng Chó nâu thì cố nhịn, nhưng rồi cuối cùng cũng bật ra: ``Mình chịu được, mình chịu
được, mình- Ahaha!'' còn nàng Sói đen chỉ biết kêu trời: ``Tại sao tôi cũng bị dính chưởng
chứ- Ahahaha!''

Tiếng cười vang lên như sóng cuộn, kéo dài suốt hơn hai mươi phút. Đến cuối cùng, bốn
người trên thanh gỗ, kiệt sức, đồng loạt cầu xin: ``Bọn em/tôi hứa sẽ không lừa dối mọi
người nữa mà!''

\il

Bên ngoài, khi xem lại đoạn này, cả phòng bật cười. Ru-chan là người đầu tiên lên tiếng,
giọng tỏ vẻ rất thỏa mãn: ``Đáng đời mấy người! Tôi đã nói là không được làm chuyện
khả nghi như thế rồi mà!''

Pekora khoái chí cười lớn, chỉ tay về phía màn hình: ``Peko peko! Đây đúng là màn xử
án mà tôi yêu thích nhất, peko!''

``Tại sao mọi người vẫn giữ lại cảnh phim đó chứ?'' Mio-chan nhìn lại bản thân bị chúng
tôi châm chọc, thốt lên đầy giận dữ.

``Thôi nào, Mio,'' Fubuki-chan lên tiếng, động viên cô nàng: ``cảnh đó xem lại cũng thấy
vui mà, đúng chứ?''

A-san thì thở dài: ``Tôi ngạc nhiên đây chính là trò đùa của cậu đấy, Hikawa-san ạ.'' Tôi
mỉm cười: ``Nhưng mà nói vui mà, đúng chứ?''

``Cậu nói cũng phải. Chắc chỉ có \hll\ chúng ta mới bày ra trò này,'' chị thở dài, ``nhưng
tôi vẫn không đồng ý chọc ghẹo kiểu này đâu.''

``Vui cái hú gì chứ! Cái này là làm nhục bọn em thì có!'' Sói đen-chan vẫn quả quyết phản
đối.

Tôi bật cười, gãi đầu: ``Kế hoạch này của anh lúc đó đúng là hơi quá tay thật. Nhưng nhìn
lại thì\dots\ đáng nhớ nhỉ?''

Mio-chan, mặt vẫn còn đỏ, cố gắng phản bác: ``Không phải chuyện đáng nhớ đâu! Cái
này đáng xấu hổ thì có!''

Koro-chan cúi thấp đầu hơn, lí nhí: ``Đúng thế\dots\ Sao lại đưa cả đoạn này vào phim chứ\dots\ Lại còn chiếu cả đoạn tôi dấm đài nữa kìa!''

Okayu-chan thì chỉ nhún vai, giọng điệu vẫn rất bình thản: ``Cái gì vui là được. Chúng ta
đang làm phim mà, phải có kịch tính một chút chứ~''

Fubuki cười khì, lè lưỡi nói với giọng tinh nghịch: ``Chắc từ giờ tôi sẽ phải tìm cách nào
đó kín đáo hơn mới được\dots''

Lập tức, Pháp sư-chan chen vào: ``Không có lần sau nào nữa đâu! Tôi sẽ giám sát tất cả
các người!''

Cả phòng lại rộ lên tiếng cười, trong khi tôi chỉ biết lắc đầu cười trừ.
Cảnh phim tiếp tục với một hình ảnh yên bình nhưng cũng đầy tiếng cười: buổi đêm giao
thừa, khi mọi người ngồi quây quần bên nồi bánh chưng nghi ngút khói. Ánh lửa bập
bùng chiếu sáng khuôn mặt từng người, tạo nên không khí ấm áp khó quên.

\il

[\dots] Nhìn quanh, tôi thấy hầu hết mọi người làm (bánh chưng) khá ổn theo chỉ dẫn của bà.
Ngoại trừ một trường hợp nổi bật.
``Thế này đúng chưa ạ?'' giọng của Fubu-chan vang lên, với một xấp lá dong cao quá mặt
khiến ai nấy đều trố mắt. Tôi phản ứng ngay: ``Thế này thì nhiều quá! Em định làm món
lá dong chưng đấy á!?''
Tiếng cười vang lên khi Fubuki cười ngượng, vội vàng bỏ bớt.

  [\dots] Tôi nhìn sang Noel-PON, thấy em ấy cầm một tảng thịt to ngang lưng lên, làm cho chúng
tôi đổ mồ hôi hột. ``Cháu có thể cho cả tảng thịt to như thế này được không ạ?''

Ru-chan phản ứng: ``Tham quá đó, Noel'' Mel-chan thêm vào: ``Nếu thế thì sao mà gói
được đây?''

Noel-PON chỉ cười trừ, nhưng ánh mắt có vẻ vẫn còn luyến tiếc tảng thịt ``hoành tráng''
kia.

  [\dots] Phù thủy-chan giơ lên chiếc bánh chưng hình con chim hồng hạc được gấp như origami.

Tôi nhướng mày ngạc nhiên: ``Em gói nó lại kiểu gì vậy?''

Mẹ tôi xen vào, tỏ vẻ nghiêm nghị: ``Các cháu gói như thế thì cogn gì là bánh chưng nữa!''
Choco-sensei đắc ý kết luận: ``Như Kuon-mama-sama nói đó. Bánh chưng phải vuông
vắn mới đúng ý nghĩa.''

Mel-chan cũng gật gù đồng tình: ``Đúng vậy, đó là biểu tượng của đất mà. Đừng có phá
hỏng hay méo mó tinh thần của nó chứ.''

[\dots] Okayu-chan nhìn chiếc bánh chưng buộc xong, nhìn hình dáng của nó có vẻ quen thuộc:
``Nhìn giống bảng chơi cờ ca-rô ghê\dots'' Thấy thế, Korone-chan phấn khích: ``Ta thử chơi
trên đó đi!''

Cặp Chó-Mèo liền tận dụng nguyên liệu thừa để ``vẽ'' bảng caro trên một chiếc bánh và
bắt đầu chơi ngay tại chỗ. Kết quả, nàng Miêu tím dễ dàng đánh bại bạn mình chỉ sau bốn
nước.

``Tôi thắng rồi meo~'' nàng Mèo thốt lên đầy thích thú. Nàng Khuyển trố mắt ra, tỏ vẻ
không phục: ``Ể~? Ogayu, không công bằng! Chơi lại!''

Pekora ngồi cạnh đó, liền lên tiếng nhắc nhở: ``Đồ ăn không phải là đồ chơi đâu, peko!
Nghiêm túc giùm đi cái!''

[\dots] Và, đoán xem người duy nhất không được tham gia vào công việc gói bánh chưng này
phản ứng gì nào? Haachama-chan thò đầu ra, hét lớn dù vẫn đang dỗi vì bị cấm cửa: ``\textbf{Các
  người sẽ hối hận vì không cho Haachama này tham gia!}''

Đáp lại chỉ là những tràng cười của tất cả chúng tôi, khiến cả một đoạn đường vốn im lìm
giờ đây tràn ngập niềm vui và sự náo nhiệt của gia đình.

\il

Roboco-chan chỉ tay vào chiếc nồi bánh chưng trên màn hình, tay lúng túng chỉ vào chiếc
bánh chưng to bất thường do chính em ấy gói: ``Ơ\dots\ Tôi lỡ tay làm nó hơi to quá. Nhưng
vẫn ăn được mà, đúng không?''

Miko-chi nghiêng đầu, đáp: ``Bánh của Roboco-san trông to lắm nhỉ? Làm kiểu gì mà ra
được cái khối đó thế?''

PON cười khì, đáp lại với vẻ ngượng ngùng: ``Tôi chỉ nghĩ là thêm nhân vào cho đầy đặn
thôi\dots\ không ngờ nó lại phình ra như vậy.''

Bên cạnh hai người họ đang bàn tán về chiếc bánh chưng khổng lồ đó, Haato -chan khoanh
tay, mặt tỏ vẻ ấm ức: ``Tôi vẫn không thể tin là mọi người không để tôi tham gia vào cảnh
này. Tại sao mọi người lại không cho tôi tham gia chứ!''

A-san, ngồi bên cạnh tôi, bật cười nhẹ: ``Cũng phải thôi, bởi vì nhiều bánh chưng để làm
quá mà. Nhưng nghĩ lại thì\dots\ có lẽ vậy lại hay hơn, nhỉ?''

Tôi mỉm cười, gật đầu, mắt vẫn dán vào màn hình: ``Đúng vậy á chị. Ai đó mà tham gia
cảnh này thì chắc nồi bánh chưng đã bay về cùng các Táo mất rồi.''

Cả phòng bật cười rộ lên. Haato-chan quay sang tôi, mắt đầy lửa giận: ``Anh nói thế là ý
gì hả Kuon-senpai!? Em cũng muốn góp sức với mọi người chứ bộ!''

Người máy-chan cười hiền, vỗ vai em ấy: ``Thì bọn chị đã cho em một vai trông nồi bánh
chưng mà. Em cũng góp một chút vào công việc này mà, đúng không nào?''

Vu nữ-chan gật gù, góp lời: ``Đúng thế ne. Làm bánh chưng mà như làm bánh khổng lồ
thì khó ăn lắm. Nhưng\dots\ nhìn cảnh này, tôi nhớ lại không khí hôm đó. Thật sự rất vui
đó, ne.''

Tôi nhìn mọi người trong phòng, cảm nhận được không khí thân thuộc lan tỏa.

Sau cảnh nồi bánh chưng, bộ phim chuyển tới một khoảnh khắc đầy sắc màu và âm thanh:
đêm pháo hoa. Ánh sáng từ pháo bừng lên rực rỡ trên màn hình, làm sáng bừng cả không
gian.

\il

[\dots] Fubuki-chan, với tinh thần ``hay chơi lớn'', en ấy mang ra một quả pháo hoa khổng lồ có
đường kính bằng một chiếc bàn học, kèm theo hàng chục quả lớn nhỏ chất đống xung
quanh.

Mio-chan nhìn cảnh tượng này, không giấu nổi vẻ hoảng hốt: ``Bà lấy mấy thứ này ở đâu
vậy!?'' Bạn thân của em ấy cười tươi rói: ``Ở kho vũ khí bí mật của Pekora-chan đấy!''

Peko-chan, đang đứng cách đó không xa, nghe thấy vậy liền giật mình: ``Này! Đống đó
là của tôi mà! Ai cho chị động vào hả, peko!?''

Nhưng chưa kịp giằng lại, Matsuri-chan, với sự năng nổ thường trực, đã nhanh tay châm
lửa vào dây pháo mà chẳng thèm nghĩ ngợi. Nàng Thỏ nhảy bổ tới, cố gắng thổi tắt dây
cháy, nhưng mọi nỗ lực đều thất bại.

``\textbf{Nguy rồi, peko! Nguy rồi, peko!!}'' em hét lớn, hoảng loạn. Thế rồi, trong một pha xử
lý ``đi vào lòng đất'', Pekora-chan quyết định ôm chặt quả pháo khổng lồ như thể muốn
hy sinh để cứu lấy mọi người.

Sự việc diễn ra quá nhanh khiến mọi người không kịp căn ngăn, để rồi khi dây cháy hết,
quả pháo đó bay lên không trung với tốc độ cao.

``\textbf{Pekora(-chan)!!!}'' tất cả chúng tôi la lên thất thanh đầy lo lắng khi thấy quả pháo bùng
cháy, bay vút lên không trung mang theo nàng Thỏ xanh xấu số.

``\textbf{Cứu tôi với!! Peko~!!}'' tiếng của em ấy vọng xuống khi quả pháo bắt đầu tỏa sáng rực
rỡ giữa bầu trời đêm.

\textbf{*BÙM!*}

Quả pháo nổ tung trên cao, tạo ra những dải sáng lấp lánh, theo sau là hàng loạt chùm
pháo hoa nhỏ bắn liên tiếp. Bầu trời đêm trở thành một bức tranh đầy màu sắc.

\il

Thỏ-chan đập bàn cười lớn đầy sảng khoái nhưng vẫn không giấu nổi sự tức giận: ``Tôi
vẫn còn cay chuyện này đấy! Cây pháo đẹp nhất của tôi bị lấy mất! Pe \uparrow ko \downarrow pe \uparrow ko \downarrow!''

Fubu-chan cố nhịn cười nhưng đôi tai cáo lại run run: ``Này, lúc đó chị mà không lấy thì
ai biết được em có dám đốt hay không? Đằng nào thì nó cũng đã cháy rất đẹp rồi mà.
Đúng không, mọi người?''

Miko-chi gật đầu, tỏ vẻ đồng tình: ``Ừ nhỉ, pháo của Pekora-chan đẹp thật. Nhưng mà,
nếu không có Fubu-chan, có khi chúng ta còn không thấy được cảnh đó á, ne.''

Okayu-chan bất ngờ lên tiếng với giọng điệu thản nhiên nhưng đầy ẩn ý: ``Pekora-chan,
chị nghĩ pháo hoa của em thật đấy. Vui là chính mà, na?''

Korone ngay lập tức bổ sung, bật cười khúc khích: ``Phải đấy! Nếu Pekora-chan cay quá
thì lần sau tụi chị sẽ chuẩn bị cho em một thùng pháo riêng luôn, được không?''

Tôi không nhịn được cười, liền góp lời: ``Nếu cần, anh sẽ tự mua cho em vài cây pháo
riêng, Pekora-chan.''

Nàng Thỏ quay sang tôi, giọng đầy nghiêm túc nhưng lại không giấu được vẻ trẻ con trong
mắt: ``Anh nói thế thì nhớ đấy! Lần sau phải chuẩn bị cho tôi nhiều vào, peko!''

Cả phòng bật cười lớn với yêu cầu có phần khá buồn cười từ em ấy.

Đến buổi sáng mùng một Tết. Ánh bình minh dịu dàng tràn ngập khung cảnh, những tia
nắng đầu tiên len qua từng tán cây và mái ngói.

Tôi vẫn không quên được buổi sáng mùng Một Tết hôm đó, khi Ayame-chan bỗng chốc
gọi tôi dậy rồi lôi tôi lên hẳn tầng thượng chỉ để chiêm ngưỡng buổi sáng đầu năm đầy
hùng vĩ này.

Máy quay nghiêng ngả, lắc lư theo từng bước chạy của em ấy. Khi cả nhóm đặt chân lên
tầng thượng, bầu trời vẫn còn lác đác vài đám mây che khuất mặt trời.

\il

[\dots] Họ cùng nhau bắt đầu hướng mắt lên trời, chờ đợi những tia mặt trời đầu tiên trong ngày
mùng Một Tết.

\dots Và đúng như dự đoán của tôi, trời vẫn u ám, mây đen dày đặc che kín, không một tia
sáng nào lọt qua nổi.

``Rõ ràng thế này mà vẫn bắt anh dậy\dots\ Mệt thật sự,'' tôi ngáp một hơi dài, lẩm bẩm ca
thán. Tôi biết rõ trời buổi sáng nay sẽ âm u rồi, nên tôi không ngạc nhiên lắm.

Nhưng, một vài người trong các cô gái lại tỏ vẻ hoang mang hơn hẳn. Nhất là Matsurichan:
``Lạ nhỉ, giờ này rồi mà tại sao mặt trời vẫn chưa lên?''

A-san đứng cạnh liền giải thích với vẻ điềm tĩnh quen thuộc: ``Tôi nghĩ là do trời có quá
nhiều mây, Natsuiro-san.''

Mio-chan thở dài: ``Thế mọi người không ai xem dự báo thời tiết à?''

[\dots] Trong lúc mọi người còn đang rối bời, đột nhiên, Pekora-chan reo lên: ``Sao chúng ta
không nhờ Miko-chan nhỉ, peko?''

Tất cả mọi ánh mắt lập tức đổ dồn về phía Vu nữ-chan, khiến chính em ấy cũng ngơ ngác:
``T-Tôi á, ne!?''

Peko-chan tiếp lời: ``Bà có cái quạt gì đó đúng không, peko?'' Roboco-chan cũng lập tức
hùa vào: ``Đúng đó, mình nhớ cậu từng dùng nó để thực hiện màn ``hoa vũ'' mà.''

``N-Nhưng cái đó thì giúp được gì, ne!?'' em ấy hoang mang, không hiểu họ đang ám chỉ
cô với điều gì.

  [\dots] Nhưng rồi, dưới áp lực của cả nhóm, ai nấy đều đồng thanh hối thúc: ``Làm đi!'' liên tục,
Miko-chan cuối cùng cũng phải hét lên: ``\textbf{Thôi được rồi!}''

Em ấy bước ra mép tòa nhà, tay cầm chiếc quạt, rồi bắt đầu thực hiện màn nhảy gọi Thần
- hay ít nhất, mọi người vẫn gọi vậy. Tôi đứng đó, vừa ngạc nhiên vừa cố nhịn cười, vừa
thầm nghĩ: ``\textit{Mình không nghĩ em ấy sẽ làm được đâu. Ý tưởng hay đấy, nhưng chắc phải
  cỡ quạt Ba Tiêu mới làm được.}''

Trong lúc nhảy, tôi còn nghe thấy Miko-chan còn lầm rầm như đang đọc thần chú: ``Kamisama,
xin hãy nghe lời thỉnh cầu của bọn con!''

Kết thúc điệu nhảy, cả nhóm im lặng chờ đợi, ánh mắt đổ dồn lên bầu trời. Không có gì
thay đổi cả.

``Đã bảo rồi mà, thôi, đừng tốn công phí sức nữa. Mau đưa anh về đi-''

Nhưng rồi, như một phép màu, gió mạnh thổi qua, xua tan đám mây đen dày đặc, để lộ
một khoảng trời quang đãng ngay vị trí mà mặt trời chuẩn bị ló dạng.

``C-Cái quái!?'' tôi há hốc mồm trong sự ngỡ ngàng. Trong khi đó, các cô gái trố mắt nhìn
cảnh tượng kỳ diệu trước mặt, rồi cùng nhau vỗ tay rầm rộ, miệng không ngừng khen
ngợi.

Haato-chan reo lớn, ôm chầm lấy tiền bối của mình: ``Chị làm được rồi đó, Miko-chan!''
Trong khi đó, đến chính người đã thực hiện nghi thức đó vẫn còn đang bối rối: ``C-Chỉ là
hên (may mắn) thôi mà, ne!''

Tôi vẫn chưa tin được Miko-chan có thể làm được như thế.

\il

Aqua-chan phá lên cười trong phòng chiếu sau khi xem lại điệu nhảy ngố tàu của Mikochi:
``Miko-senpai múa điệu thần linh trông buồn cười quá!''

Miko - người đang xem lại đoạn phim - đỏ mặt, quay sang em ấy phản bác: ``E-Em không
cần phải nhắc lại đâu! Chẳng qua là\dots\ lúc đó chị hơi quá nhập tâm thôi mà!''

Choco-sensei cười mỉm: ``Nhưng dù gì thì Miko-sama đã giúp chúng ta ngắm được buổi
bình minh năm mới mà, đúng chứ?''

Aki nở một nụ cười mãn nguyện: ``Dù gì thì chúng ta cũng có được một buổi sáng đầu
năm đầy may mắn rồi.''

Cảnh quay kết thúc bằng tiếng cười rộn ràng của cả nhóm trên tầng thượng, hòa lẫn với
âm thanh của một buổi sáng Tết bình yên.

Tôi mỉm cười, quay sang cả nhóm trong phòng: ``Chính những khoảnh khắc như thế này
đã làm nên một cái Tết khó quên, đúng không?''

Mọi người gật đầu đồng ý, không ai nói thêm lời nào, nhưng nụ cười trên môi họ đã thay
lời muốn nói.

Cảnh tiếp theo trong bộ phim chuyển tới một không khí sôi động tại sân nhà, nơi cuộc thi
giã bánh mochi diễn ra. Âm thanh của chày gỗ đập xuống cối vang vọng, xen lẫn tiếng
cười nói rộn ràng.

\il

[\dots] Ayame nhảy lên một bậc thềm gần đó, hô to: ``Thông báo! Cuộc thi giã bánh mochi chính
thức bắt đầu! Ai giã nhanh và đẹp nhất sẽ được thưởng\dots\ một phong lì xì từ Kaien-san!''

Mẹ đang bận lau sạch đôi dép bị dính nước, giật mình ngẩng lên: ``Ơ từ từ, cô chưa chuẩn
bị tiền lì xì-'' nhưng chưa kịp dứt lời, nàng Quỷ sừng tinh nghịch đã bồi thêm một câu đầy
sát thương: ``Và phần thưởng đặc biệt\dots\ là một đêm ngủ với Kuon-senpai!''

[\dots] Cả sân như đông cứng lại trong vài giây. Các cô gái quay đầu nhìn Ayame, rồi bỗng dưng
mặt của họ đỏ hỏn hơn hẳn.

``M-Một đêm ngủ chung với Kuon-senpai sao?'' tất cả bọn họ đồng thanh.

\il

Cả nhóm trong phòng chiếu cùng bật cười. Pekora-chan cười lớn, chỉ tay vào màn hình:
``Ayame này đúng là lắm trò!''

A-san lắc đầu, nhưng khóe miệng lại cong lên thành nụ cười: ``Quả thật đúng là như vậy
đấy\dots''

Ru-chan ở ngoài đời vẫn đỏ mặt, thì thầm đủ để mọi người nghe thấy: ``N-Nhưng mà\dots\ em cũng đã thắng rồi\dots''

``Nè, mình cũng là người chiến thắng nữa mà!'' Noel-chan cụng vào vai của cô nàng, cười
méo xệch.

Cả phòng lại bùng nổ tiếng cười, nhất là khi trên màn hình, cảnh cặp đôi Chó-Mèo bị Ruchan
đạp đổ tận hai lần, Ayame và Noel cố gắng phối hợp nhau nhưng rồi bột rơi vương
vãi. Những khoảnh khắc hỗn loạn ấy hiện lên rõ nét, khiến cả nhóm không thể ngừng
cười.

Kết thúc cảnh quay, tôi mỉm cười nhìn mọi người: ``Cuộc thi này đúng là đáng nhớ, phải
không?''

Pekora cười lớn, giơ tay lên: ``Chắc chắn rồi! Em vẫn cảm thấy lâng lân đến tận bây giờ
đây, peko!''

Cảnh tiếp theo trong bộ phim là cuộc hành trình trở về quê nội đầy kịch tính, mà tôi vẫn
không thể tin được là nó lại thực sự xảy ra trên đường phố Etsunan.

\il

[\dots] Hai chiếc xe do Fubuki và Okayu lái phóng như tên lửa trên con đường quốc lộ. Cảnh
tượng không khác gì một màn chơi A(s)phalt sống động: tiếng động cơ gầm rú, bánh xe
nghiến trên mặt đường, và cả những tiếng la hét từ hành khách phía sau.

Fubuki, mắt sáng rực như đã vào ``chế độ game thủ'': ``Cứ thế này tôi sẽ về nhất cho xem!''
Đối thủ của cô, Okayu, tay vững vàng trên vô lăng, bật cười: ``Không, không, là tôi chứ~''

Mio cố giữ không ngã khỏi ghế sau, hét lớn: ``\textbf{Đừng có tự biến mình thành hung thần
  đường phố chứ!!}''

Korone phấn khích, đến mức không ngừng cổ vũ: ``Ogayu cố lên, cố lên nào~'' dường
như ngó lơ những hành khách đang tái mét mặt mày ở phía sau.

  [\dots] Đúng lúc đó, Fubuki trông thấy một chiếc xe tải lớn đang chạy phía trước, trên đó có một
ván dốc chở hàng. Mắt cô sáng lên: ``Nhìn nè, Mio! Một cú xoay vòng tuyệt hảo đang
chờ chúng ta!''

Noel thấy vậy, nổi hứng hơn hẳn: ``Tới luôn đi, Fubuki-senpai!!'' mặc cho những lời la ó
can ngăn của những người còn lại trên xe.

Nàng Cáo trắng đánh lái và tăng tốc, lao thăng lên ván dốc xe tải. hiếc xe lập tức nhấc
bổng lên không trung, xoay một vòng tròn hoàn hảo trên không\dots

\begin{dialogue}
  \speak{Miko} Cái gì vậy ne\dots?
  \speak{Mio} \textbf{Fubuki!!}
  \speak{Fubuki} Woo-hoo!
  \speak{Noel} Ta đang xoay vòng trên cao!
\end{dialogue}

\dots trước khi đáp xuống mặt đường với một cú nảy mạnh, khiến những người đi đường
thêm một phen khiếp đảm.

Mio lúc này tái mét, cả người cô mềm nhũn ra: ``Tôi sẽ không đi chung xe với bà một lần
nào nữa đâu!'' Rushia và Miko ôm chặt lấy nhau, gần như khóc:

\begin{dialogue}
  \speak{Miko} T-Thật là đáng sợ quá, ne\dots\ Chị tè ra quần mất thôi!
  \speak{Rushia} B-Bình tĩnh lại đi nào, Miko-chan! Fubuki-senpai, dừng lại đi!
\end{dialogue}

Ở phía sau, xe của Okayu cũng không chịu kém cạnh. Korone, mắt sáng rực như đèn pha,
chỉ tay về phía một chiếc ô tô đỗ bên lề đường: ``Ogayu! Chúng ta phải làm gì đó thật ấn
tượng hơn!''

Nàng Miêu liếc nhanh, nhoẻn miệng: ``Bà đang có bài gì sao, Koro-san?'' Nàng Khuyển
hét lớn, giơ nắm đấm lên cao không tỏ vẻ do dự gì: ``\textbf{Lộn vòng đi nào!}''

A-chan vội lên tiếng: ``N-Này! Các cô tính làm cái trò gì vậy!?'' nhưng dường như câu
nói của nàng Đeo kính không thể đến tai của cặp đôi ``quỷ sứ'' này.

``Vậy thì lên nào!'' Okayu cười, đánh mạnh vô lăng. Xe lập tức xoay một vòng tròn trên
đường nhựa, trước khi bất ngờ đâm nhẹ vào mép vỉa hè, nhấc bánh sau lên, và\dots\ chiếc xe
lao vào không trung như một pha nhào lộn trong phim hành động của xứ cờ hoa.

Những hành khách trên chiếc xe này cũng la hét lớn không kém gì xe đi đằng trước họ:

\begin{dialogue}
  \speak{A-chan} Inugami-san! Nekomata-san! Các cô sẽ bị trừ lương vì chuyện này\dots!
  \speak{Aki} Eeek!
  \speak{Choco} Hai em đang làm cái trò gì thế hả!?
  \speak{Haato} \textbf{Xoay vòng nào, cặp đôi \textit{GAMERS} ơi!}
\end{dialogue}

Chiếc xe đáp xuống với một cú nảy mạnh, tông hạ luôn chiếc ô tô đỗ bên lề, khiến nó bật
tung ra một góc.

  [\dots] nàng Dị giới bất ngờ cất giọng lo lắng: ``Ờm, mọi người\dots?'' Tất cả mọi người
ở cả hai xe đồng loạt vang lên: ``Sao!? / Cái gì!?''

``Phía trước có tàu hỏa kìa\dots''

Tất cả đều quay đầu nhìn về phía trước. Một đoàn tàu hỏa khổng lồ đang chậm rãi tiến
lại gần, và hàng rào chắn đang từ từ hạ xuống.

``Mau dừng xe đi chứ, Okayu-sama! Đừng có nghĩ đến chuyện dại dột!'' nàng Y tá hoảng
hốt. Nàng Pháp sư ở xe của Fubuki cũng lo lắng không kém: ``Em là Cô đồng nhưng em
cũng không muốn chết đâu!!''

Trong khoảnh khắc đáng lẽ phải dừng lại, những ``quái xễ'' lại đồng loạt nảy ra một ý
tưởng ``táo bạo''. Dường như họ chẳng đếm xỉa gì tới những tiếng la hét thảm thiết của
các hành khách phía sau.

``\textbf{Ta cùng vượt đường ray thôi!!}''

Noel và Roboco, cặp đôi ngốc nghếch nhưng hăng máu (vì không sợ chết) nhất, reo hò:
``Vượt rào! Vượt rào!'' Haato thậm chí còn hét lớn: ``\textbf{Nhanh nữa lên nào!!}''

Miko lúc này đã sợ đến mức ngất xỉu, khiến cho Rushia ngồi cạnh: ``Miko-chan!? Mikochan!!
Đừng chết mà!'' Nàng Vịt cũng không kém cạnh: ``Đây có phải là GTA đâu!!''

Nàng Song tính tiếp tục hú hét, cổ vũ cho cả hai người: ``\textbf{Cố lên nào mọi người!}'' Lập
tức, Subaru không kiềm được cơn giận: ``Lại còn cả chị nữa!! Ngậm miệng lại cho em
đi!!''

Hai chiếc xe lao tới với tốc độ không tưởng, bất chấp hàng rào chắn tàu hỏa đang hạ xuống
hoàn toàn. Khi đoàn tàu chỉ còn cách vài mét, cả hai xe đồng loạt tông sầm vào hàng rào
chắn, làm nó bật tung lên.

Toàn bộ những người bên trong xe lại tiếp tục la hét. ``\textbf{Chúng ta sẽ chết mất! Chúng ta
  sẽ chết mất thôi!!}'' một vài người trong số họ thất thanh, trong khi một số người thì kích
động vui sướng tiếp tục cổ vũ.

Trong tích tắc, hai xe vượt qua đường ray ngay trước khi đoàn tàu lao qua. Tiếng còi tàu
rú vang, làm chấn động cả khu vực. Những người tham gia giao thông nhìn cảnh tượng
đó trố mắt ra ngỡ ngàng, không thể tin vào những gì họ đang được chứng kiến.

  [\dots] ia đình Kuon đang lái xe một cách thong thả trên con đường làng. Tiếng gió vi vu hòa
cùng cảnh đồng quê yên bình khiến không khí vô cùng dễ chịu. Cả bốn người bọn họ vừa
đi, vừa cảm nhận được khí trời đang thổi vào người họ, vừa trò chuyện với nhau vui vẻ
giống như trong gia đình bình thường.

Bỗng nhiên, họ nhìn thấy năm khối lập phương có dấu hỏi đang xoay vòng, lơ lửng giữa
đường. Ryujin nhíu mày nhìn những khối lập phương kỳ lạ: ``{\color{red} Cái gì thế này? Đồ để trang trí xe máy à?}''

Ryuuhou ở xe bên kia, hớn hở nói: ``{\color{red}Cái này con không chắc\dots\ Nhưng chúng trông giống
  power-up trong Ma**o Kart đấy bố!}''

Kuon nghiêng đầu ngạc nhiên: ``{\color{red}Em nói như thể đó là trò chơi vậy\dots\ Mà tại sao nó lại ở
  đây được chứ? Đây là Kawachi, chứ có phải thành phố Điện tử-}''

Tuy nhiên, vì mải quan sát, Ryujin đã vô tình đâm thẳng vào một trong những khối lập
phương làm nó vỡ tung ra. Ngay lập tức, một vòng quay xuất hiện bên cạnh họ, xoay tít
như đang chọn giải thưởng.

``{\color{red}Cái quái gì vừa diễn ra vậy?}'' người bố ngạc nhiên khi vòng quay đang xoay một cách
chóng mặt. Chưa để ai nói gì thêm, vòng quay dừng lại, hiện lên biểu tượng hai mũi tên
màu cam chỉ hướng lên trên.

Ông chớp mắt khó hiểu: ``{\color{red}Đây là biểu tượng gì đây?}'' Cậu Tổng ngồi đằng sau, cầm lấy
hai mũi tên trên và gắn thử vào chiếc yên xe máy ở thùng sau xe: ``{\color{red}Để con thử xem nào\dots}''

Ngay khi mũi tên vừa chạm vào yên xe, chiếc xe máy rít lên một tiếng và\dots\ phóng vọt
lên trời như tên lửa! Người mẹ cầm lái trên xe ở bên cạnh, hoảng hốt hét lên: ``{\color{red}\textbf{Chuyện
    gì thế này!? Xe biết bay ư!?}}''

Cô em gái ngửa mặt lên trời, chỉ vào chiếc xe của hai bố con đang bay cao: ``{\color{red}Onii-san
  đang bay lên trời luôn rồi! Trông như siêu nhân vậy!}''

Kuon dù đang ôm chặt bố, vẫn có thể vọng lại: ``{\color{red}Đừng có cười anh chứ!}'' Đến Ryujin
cũng hoang mang không kém: ``{\color{red}Chúng ta tại sao lại ở trên trời vậy!?}''

\il

Mio-chan xem đến đoạn này, không khỏi ôm đầu: ``Em thề lúc đó tim em như muốn rớt
ra ngoài luôn! Có khi thành phố của anh bị nhiễm sự quấy phá từ mọi người cũng nên á,
Kuon-senpai!!''

Tôi chỉ biết cười khổ, nhún vai: ``Anh vẫn chưa tin là chúng ta có cả Aspha*t và Mar*o
Kart trên đường phố đấy\dots''

Aki-chan, người điều khiển xe của Okayu trong phim, thở dài ôm trán: ``Nghĩ lai thì mình
vẫn còn thấy hoảng\dots\ Thế quái nào mà xe của em ấy có thể bay lên được nhỉ?''

Okayu-chan chỉ cười, thản nhiên nói: ``Thế thì mới gọi là Lamb*rghini xịn chứ~''

Subaru-chan phấn khích, vung tay, hoàn toàn ngược lại với vẻ hoảng loạn trong phim:
``Cái đoạn vượt tàu đó ngầu thật sự! Tôi mà lái chắc không làm nổi đâu!''

Haato-chan chăm chú màn hình, mắt sáng lên: ``Cái này giống hệt phim hành động thật
này! Ước gì tôi được làm người lái!''

Mel thì tỏ vẻ e thẹn hơn với cảnh đó: ``May mà chị không ở trong mấy chiếc xe đó\dots\
Nhưng nhìn thôi chắc chị cũng muốn xỉu rồi.''

Tôi lắc đầu, quay sang A-san, thầm thì: ``Nhìn lại chuyện này, em vẫn không tin là mọi
người sống sót được sau ngày hôm đó.''

A-san mỉm cười: ``Cậu cũng thế mà. Chắc đó là nhờ may mắn của cậu đấy, Hikawa-san.
Nhưng đừng thử lại lần nữa nhé.''

``Đến bản thân em cũng sợ vì mấy tai nạn hài hước thế này rồi chị ạ\dots''

Bộ phim chuyển sang không khí truyền thống hơn, với cảnh mọi người tham gia hoạt
động chúc Tết và mừng tuổi vào mùng Một và mùng Hai.

Trong phim, mùng Một là ngày chúng tôi dành cho họ hàng bên nội. Đại diện bốn người
từ \hll\ - Sora-chan, Mel-chan, Fubuki-chan và Mio-chan - cùng nhau đi thăm từng
nhà, từ các cụ cao tuổi đến những gia đình trẻ hơn. Họ ngạc nhiên khi thấy một nhóm
đông đúc như vậy, nhưng cũng nhiệt tình đón tiếp, đặc biệt khi biết đây là các ``khách quý
từ Nhật Bản.''

\il

[\dots] Ngôi nhà của một người họ hàng nhà nội Kuon nằm sâu trong một con ngõ nhỏ, bao quanh
bởi những ngôi nhà mặt tiền to lớn. Khi chiếc xe vừa đỗ lại, cả nhóm nhanh chóng xuống
xe, chỉnh trang trang phục trước khi bước vào bên trong ngõ.

Bên trong con ngõ hơi tối tăm nhưng đã được dọn dẹp sạch sẽ đó là đầy rẫy những ngôi
nhà khác. Nhà của một người họ hàng có thể phân biệt được bằng hai lớp cửa: một chiếc
cửa kéo tay bằng nhôm ở bên ngoài, và một cánh cửa kính đẩy ở bên trong.

Cả hai bố con mở cánh cửa kéo ra, rồi cùng nhau đẩy chiếc cửa kính vào, lộ ra hai người
đang ngồi trên ghế bành trắng đang chờ đợi bọn họ: một người đàn bà lớn tuổi, tóc đã
điêm những chấm bạc, và một người đàn ông trông trẻ hơn Ryujin, đầu gần như trọc lốc.
Người đàn ông đó cũng là người làm việc lâu năm với ông đã được gần mười năm.

``{\color{red} Bà, bọn cháu đến chúc Tết đây ạ,}'' người bố nói, theo sau là cậu Tổng chào hai người
bằng phép lịch sự tối thiểu: ``{\color{red} Cháu chào bà, cháu chào chú Min(h).}''

Các cô gái cũng đồng loạt cúi đầu chào, khiến cả hai người họ bất ngờ. Người phụ nữ lớn
tuổi nhìn hai bố con cậu, rồi quay sang các cô gái, ánh mắt lộ vẻ tò mò: ``{\color{red} À, đây chẳng
  phải là bạn của thằng Cò sao?}''

Ryujin gật đầu, giải thích: ``{\color{red} Dạ, đây là Mel-chan, Sora-chan, Mio-chan và Fubuki-chan
  - đồng nghiệp làm việc chung với Kuon. Các cháu lần đầu đến Etsunan nên con mời về
  chơi dịp Tết.}''

Người đàn ông nhếch mép cười, nhìn cậu Tổng rồi giở giọng trêu: ``{\color{red} Thằng Kuon này cũng
  gớm phết đấy! Được làm việc với gái như thế sướng quá còn gì!}''

Câu đùa của người đàn ông khiến cậu con trai đỏ mặt, còn các cô gái thì cười rúc rích.

  [\dots] ến khi kết thúc lượt chúc Tết ở nhà anh trai của Hijaki - tức anh trai trưởng của gia đình
Kuon, cô nàng than vãn: ``Có từng này thôi à\dots''

Mio nhíu mày nhìn cô bạn: ``Lì xì không phải là khoản thu nhập đâu, Fubuki. Còn chưa
kể bà phải mừng tuổi lại bọn trẻ con đấy.''

Nàng Dơi tỏ vẻ đồng cảm với cô, rồi nói với vẻ hãnh diện: ``Giống như chúng tôi làm hồi
sáng nay chẳng hạn.''

Sora mỉm cười: ``Đúng rồi. Quan trọng là cái lộc từ người ta trao cho mình đó, Fubukichan
à.''

Dù vậy, nàng Cáo trắng vẫn ca thán: ``Nhưng lộc như thế này thì ít quá! Làm sao mà đủ
để gacha đây\dots''

Kuon chỉ biết lắc đầu, bật cười nhẹ: ``Em cứ như bọn trẻ con bây giờ ấy. Thôi, sang nhà
tiếp theo đi!''

\il

Fubu-chan trong đời thật, vừa xem vừa cười híp mắt, nói với vẻ háo hức: ``Không biết tôi
được bao nhiêu tiền lì xì nhỉ~?''

Mio-chan nghe vậy, lập tức nghiêm chỉnh lại: ``Đừng có tham thế chứ\dots''

Sora-chan, dẫu là ở trong phim hay ngoài đời, luôn là hình mẫu lịch sự và duyên dáng.
Cô nàng dẫn đầu các màn chào hỏi, làm cả gia đình tôi ấn tượng.

Mel-chan thì điềm đạm hơn, nhưng cũng không quên làm tan bầu không khí với những
câu chuyện hài hước: ``Hình như em có nhớ đã đừng kể cho họ nghe về Tết Nhật rồi nhỉ?''

Đến mùng Hai, bầu không khí có phần căng thẳng hơn, vì đó là ngày dành cho họ hàng
bên ngoại, những người mà - tôi phải nhắc lại lần nữa - đa phần tôi không hề biết mặt,
cũng không hề có nhu cầu biết mặt họ. Dĩ nhiên, tôi cũng không thoát được những câu
hỏi khó đỡ từ mấy người lớn tò mò thích săm soi đời tư tôi.

\il

[\dots] Tuy nhiên, mọi thứ không phải lúc nào cũng suôn sẻ. Khi vào đến nhà thứ hai, một vài
người họ hàng của tôi - mà thú thực tôi chẳng nhớ tên hay mặt mũi ai, cũng chằng phải là
điều mà tôi bận tâm tới - bắt đầu đặt ra những câu hỏi mà chỉ nghe thôi đã muốn thở dài.

Một người phụ nữ trung niên, có lẽ là một trong các dì (về lý thuyết) của tôi, cười giả lả
hỏi: ``{\color{red} Cháu Kuon, năm nay đã có người yêu chưa? Có bạn gái chưa?}''

Tôi nheo mắt nhìn bà ấy, cố gắng giữ nụ cười xã giao: ``{\color{red} Cháu đang bận học và làm việc,
  cũng chưa nghĩ đến chuyện đó ạ.}'' Dù vậy, bà ấy vẫn không buông tha: ``{\color{red} Thế á? Nói như
  cháu á, chắc chẳng ai thèm đâu đấy hở? Kén chọn quá đó!}''

Đó THỰC SỰ là câu trả lời làm tôi muốn tức cái đầu. Tuy nhiên, ngay khi tôi định đáp
lại thì Mio-chan bước tới, mỉm cười dịu dàng nhưng lời nói thì sắc bén hơn hẳn: ``{\color{red} Dạ,
  chuyện tình cảm là chuyện riêng tư, mong dì đừng làm khó Kuon-senpai ạ.}''

``{\color{red} Với cả Kuon-senpai có bạn gái hay không cũng không phải là nghĩa vụ mà dì phải biết
  đâu ạ,}'' Fubuki-chan tiếp lời.

Người phụ nữ kia thoáng đỏ mặt, nhưng không nói thêm gì.

\ct

Sau đó, một người đàn ông khác - có lẽ là một trong các bác - lại chuyển sang hỏi về
chuyện học tập. Đó không phải là chủ đề mà tôi thấy khó chịu gì, nhưng năm nào bố mẹ
tôi cứ khen tôi học giỏi lắm, nên gần như mọi người toàn đặt nặng kỳ vọng vào tôi cũng
đủ khiến tôi đau đầu.

Bác ấy xởi lởi hỏi: ``{\color{red} Cháu đang học đại học, đúng không? Thế bao giờ ra trường? Có
  định học lên nữa không, hay sẽ ở nhà làm việc phụ bố mẹ?}''

Tôi đáp lại một cách khá nhẹ nhàng: ``{\color{red} Cháu cũng không chắc ạ. Chắc là cháu sẽ học
  khoảng hai, ba năm nữa. Còn cháu cũng không chắc là sẽ ở nhà phụ việc bố mẹ đâu ạ-}''

Tôi đang trả lời thì Pekora-chan cắt ngang lời tôi: ``{\color{red} Vì đến bản thân anh ấy còn chưa lo
  xong mà thi làm sao lo được cho gia đình đâu ạ, peko.}''

Tôi thở dài, chỉ về phía em ấy và hoàn thành câu nói: ``{\color{red} \dots như em ấy nói đấy ạ.}'' Ông bác
ấy nghe câu trả lời của chúng tôi, bắt đầu giở giọng xỉa xói: ``{\color{red} Ui dời, cháu kém thế! Thời
  ngày xưa ấy, chú đã bắt đầu kiếm tiền được rồi!}''

Câu trả lời của bác khiến tôi không thoải mái, nhưng trước khi tôi kịp trả lời, Subaru-chan
- với phong cách thẳng thắn thường thấy - đã chen vào: ``{\color{red} Bác ơi, thời của Bác xưa như
  Diễm rồi. Việc học tập của Kuon-senpai chắc chắn là tốt rồi. Bác không cần lo lắng đâu.
  Với lại, anh ấy còn làm việc ở một công ty lớn mà!}''

Người đàn ông kia hơi khựng lại, không ngờ bị nàng Vịt ``phản pháo'' nhanh đến thế.

\ct

Chưa dừng lại ở đó, một người họ hàng khác cũng hỏi han tình hình cuộc sống của tôi.
Nhưng lần này thậm chí còn lỗ mãng hơn, khi người đó lại bắt đầu so sánh tôi với con
của họ. ``{\color{red} Con tôi năm nay mới tốt nghiệp đại học, mà đã được nhận vào một công ty lớn
  với mức lương cao. Cháu thì sao, hở?}''

``\textit{Việc con cháu nhà chú thì liên quan gì đến tôi cơ chứ\dots}'' đó là điều mà tôi muốn thốt ra,
thậm chí còn muống ``đôi co'' với người họ hàng bất lịch sự đó, nhưng trước khi kịp nói
gì, Ayame-chan - với vẻ ngoài hăng như bò tót cùng tính cách sắc sảo - lên tiếng đầy mỉa
mai: ``{\color{red} Vậy sao ạ? Thế thì chúc mừng anh chị và cháu ạ. Nhưng em nghĩ mỗi người có con
  đường riêng. Anh chị đi đường anh chị, còn Kuon-senpai sẽ đi đường của Kuon-senpai,
  đúng không ạ?}''

``{\color{red} Với cả, chú có thấy bất lịch sự không khi nói phả thẳng mặt thế ạ?}'' Shion-chan lên tiếng
bênh vực tôi, dẫu rằng bản thân em ấy cũng là một kẻ thô lỗ, làm cho tôi cảm thấy hơi sai
sai.

Dù vậy, lời nói của hai em ấy như một cú đấm thẳng mặt, khiến người họ hàng kia chỉ
biết cười gượng và ngậm miệng.

Bầu không khí có phần căng thẳng dần trở lại bình thường khi chúng tôi tiếp tục chuyến
đi. Tôi thầm cảm ơn các cô gái vì đã đứng ra đỡ lời cho tôi. Dù đôi lúc họ có hơi phiền
phức, nhưng những khoảnh khắc thế này khiến tôi nhận ra rằng, họ thực sự là một phần
quan trọng của gia đình \hll\ mà tôi gắn bó.

\dots Tôi thề, tôi sẽ KHÔNG BAO GIỜ về quê ngoại thêm một lần nào nữa.

\il

Subaru-chan xem lại đoạn này, bật cười sảng khoái: ``Trông em thật tài giỏi trong tài ứng
biến đúng không nào?''

Mio-chan chỉ mỉm cười, nhẹ nhàng chỉnh lại: ``Thật ra chị đã phải nghĩ nát óc đấy\dots\ may
là không lộ.''

Miko-chan khoanh tay lại, tỏ vẻ bất bình với những người ở trong đoạn phim đã che mặt
(bởi tôi) kia: ``Họ đúng là bất lịch sự thật đấy, ne. Người ta đã không muốn trả lời rồi lại
còn đè ra hỏi nữa!''

Tiếp theo, bộ phim chuyển đến những hoạt động vui nhộn trong ngày Tết, nơi các cô gái
thử sức với những trò dân gian truyền thống của Etsunan, như tạt lon và nhảy lò cò, cùng
với trò đánh cầu Hanetsuki của Nhật Bản.

Cảnh đầu là lúc chơi nhảy lò cò, và đương nhiên, các thành viên càng thể hiện rõ hơn sự
khác biệt về khả năng giữ thăng bằng. Tôi nhớ trong phim, Shion-chan loạng choạng đến
mức ngã nhào, còn Fubuki-chan thì lại chơi bài ``tư bản'' để lừa Akutan.

\il

[\dots] Tất cả mọi người đều đang bàn tán xôn xao. Aki-chan giơ tay, hỏi: ``Vậy, luật chơi cụ thể
là gì vậy, Kuon-senpai?''

``Luật đơn giản thôi,'' tôi giải thích. ``Mỗi người sẽ lần lượt nhảy bằng một chân qua từng
ô, không được đụng vạch kẻ. Nếu đụng vạch, phải quay lại điểm xuất phát. Nhưng ở các
ô nối với nhau, người chơi phải nhảy liền hai chân cùng lúc. Ai đi hết con đường trước sẽ
thắng.''

``Tưởng dễ mà khó phết đấy,'' Mio-chan gật gù. Korone-chan giơ hai tay lên nắm lại, thốt
lên đầy hứng thú: ``Nhưng mà trông cũng hay lắm á!''

``Đúng vậy,'' tôi nháy mắt. ``Nhưng nếu ai làm rơi đồ vật mà anh ném vào ô số nào đó,
phải chịu một hình phạt đặc biệt.''
``Hình phạt gì vậy?'' Roboco-PON hỏi, tỏ vẻ thắc mắc trong sự hiếu kỳ. Tôi mỉm cười
đầy bí ẩn, đáp gọn ghẽ: ``À, lát nữa chơi thì biết.''

Khi trò chơi bắt đầu, đúng như tôi dự đoán, có hàng loạt pha ngớ ngẩn đến mức khiến cả
sân náo nhiệt bởi những tiếng cười không dứt. Ai cũng tranh nhau nhảy, nhưng việc giữ
thăng bằng trên một chân trong khi ô thì nhỏ, đường thì cong queo, chẳng dễ chút nào.
Và tất nhiên, những hình phạt kỳ lạ do tôi nghĩ ra càng làm trò chơi trở nên ``đậm chất
\textit{hologra}'' hơn bao giờ hết.

\il

Pekora-chan cười hả hê khi nhớ lại: ``Shion-chan cũng nhảy lò cò ư!? Nhìn thôi cũng thấy
buồn cười rồi, peko!''

Shion quay sang tôi, cố gắng biện minh: ``Em không có tập trung thôi mà! Nếu làm lại
chắc chắn sẽ khác!''

Tiếp theo là chơi tạt lon, và đương nhiên, sự hỗn loạn là điều không thể tránh khỏi.

\il

[\dots] Trò này không cần dụng cụ cầu kỳ, chỉ cần một chiếc lon, vài đôi dép hay giày, và một
khoảng không gian đủ rộng là đủ. Đối với các cô gái, họ đều tròn mắt như thể tôi vừa giới
thiệu một thứ gì đó từ hành tinh khác. Cũng phải thôi, quê tôi mới chơi trò này, còn họ
thì không.

``Tạt lon? Là cái gì thế?'' Aki-chan ngơ ngác hỏi. Noel tiếp lời, tỏ vẻ băn khoăn: ``Có
phải giống như ném đồ vật vào mục tiêu không?''
Shion cũng không giấu nổi sự tò mò: ``Nếu thế thì dễ quá\dots\ nhưng nghe kỳ kỳ thế nào
nhỉ.''

Tôi bật cười trước sự ngơ ngác của mọi người. Cũng đúng thôi, trò này có lẽ chỉ phổ biến
ở quê, còn với các cô gái từ thành phố lớn hay đến từ những thế giới tưởng tượng như
Aki-chan, hẳn đây là điều mới mẻ.

``Ừ, gần đúng đấy, nhưng mà thêm chút hành động chạy trốn nữa. Để anh giải thích luật
nhé.''

[\dots] Theo lệnh của Miko-chi, tất cả mọi người đồng loạt ném giày dép của mình về phía chiếc
lon. Tiếng rầm rầm vang lên khi các vật dụng bay tới tấp. Nhưng bất ngờ, giữa cơn mưa
giày dép, chiếc chùy của Noel-chan lại lao vút lên trời, vượt qua chiếc lon, và đáp xuống
rất gần chỗ của đối phương, khiến em ấy giật nảy mình.

``\textbf{Em định giết chị sao, Noel!?}'' Miko-chan thảng thốt hét lên, hú vía sau cú ném trên.
Nàng Hiệp sĩ cười trừ: ``Xin lỗi~!''
Miko-chan vội vàng lùi lại, nét mặt vừa bực bội vừa sợ hãi. Tôi đứng từ xa mà không nhịn
được cười. Quay sang Noel, tôi thắc mắc: ``Sao không dùng giày dép của em để ném?''

``Em không cởi ra được\dots\ Dép geta này chặt quá\dots'' nàng Hiệp sĩ giơ chân lên, cố kéo
chiếc dép của mình ra nhưng không thành. Cả nhóm phá lên cười khi nhìn thấy cảnh
tượng đó.

\il

``Quả đó em đã thực sự suýt giết chị đấy, Noel!'' Miko-chi nhìn lại khuôn mặt trắng bệch của mình
khi nàng Hiệp sĩ ném cái chùy về phía em ấy thay vì chiếc dép.

Noel-chan vội vàng xua tay, mặt đỏ bừng: ``Em không cố ý thật mà!''

Mio-chan nhìn hai cô nàng đang đỏ hỏn vì ngượng ngùng tai nạn đó, nói nửa đùa nửa thật:
``Ai bảo chị vẽ ô rộng quá làm gì. Cái này gọi là nghiệp quật đấy.''

Cảnh sau cùng là lúc mọi người cùng chơi Hanetsuki, trò đánh cầu truyền thống của Nhật
Bản. Ai đánh hụt sẽ bị vẽ một nhọ nồi trên mặt.

Choco-sensei vừa xem lại trận đấu đánh cầu đầy ác liệt của chín họ, vừa lắc đầu cười nhẹ:
``Cô không nghĩ trò này chúng ta chơi lại hỗn loạn đến vậy. Nhìn ai cũng lấm lem hết cả.''

Sora-chan dịu dàng bình luận: ``Dù sao thì cũng vui mà. Ai cũng cố gắng hết sức cả rồi.''

Pekora-chan thì không bỏ lỡ cơ hội để trêu chọc mọi người: ``Nhưng mà công nhận tôi
Shion với Aqua có nhiều chiêu bài thật đấy, peko!''

Shion-chan nhướng mày về phía nàng Thỏ: ``Em với Matsuri-senpai cũng giở lắm chiêu
làm khó đối thủ đấy thôi.''

Fubuki cười lớn: ``Lúc đó chúng ta thi đấu hăng say thật đấy! Nhiều khi chị tưởng chúng
ta đang ở một trận đấu sinh tử không bằng!''

Ru-chan khẽ cười đáp nhẹ: ``Dúng là như vậy đấy-nanodesu.''

Cả nhóm cười rộ lên khi nhớ lại cảnh mà họ đã có một trận đấu đánh cầu Hanetsuki đầy quyết liệt nhưng không kém phần hỗn loạn và hài hước.

Cảnh phim chuyển sang bàn ăn, nơi mọi người đang thưởng thức các món ăn truyền thống
của Etsunan. Những món ăn như bánh chưng, bánh tét, dưa hành, và cả những đĩa nem
rán giòn tan được bày biện khắp bàn. Không khí rộn ràng nhưng cũng không kém phần
ấm cúng.

\il

[\dots] Khi bữa tối bắt đầu, bàn ăn trở nên nhộn nhịp với các món đặc sản ngày Tết của vùng
Etsunan. Noel là người nổi bật nhất, với tốc độ ăn gà luộc nhanh đến mức làm mọi người
tròn mắt.

Nàng Vịt cười lớn, giở giọng trêu: ``Noel-chan! Em ăn như thể chưa được ăn gà bao giờ
ấy!'' Nàng Hiệp sĩ không đáp, chỉ ăn món thịt gà đầy ngon lành, như thể đã làm tâm
nguyện của cô lâu năm được hoàn thiện.

Mio cắn một miếng nem rán do Matsuri và Okayu làm (được Ryuuhou hướng dẫn), gật gù
đồng tình: ``Ryuuhou-chan nêm nếm vừa miệng thật. Chị phải học cách làm mới được!''

Một nửa số thành viên thích thú khám phá món bánh chưng, trong khi nửa còn lại thì hơi
chần chừ với hương vị lạ lẫm. Korone, với tinh thần ``không ngán món gì'', mạnh dạn ăn
một miếng lớn rồi thốt lên: ``Ngon quá! Mọi người thử đi, đừng sợ!''

Nàng Mèo bên cạnh bật cười: ``Koro-san làm cứ như thể mọi người vẫn còn lạ lẫm với món đó
á\~{}''

Ở góc bàn, Sora và Mel đang gắp thử món thịt đông. Nàng Dơi vừa ăn vừa nhíu mày tò
mò: ``Món này lạ thật, nhưng cũng thú vị.'' Nàng Thần tượng mỉm cười: ``Ừ, có chút mát
lạnh, ăn rất hợp vào dịp Tết.''

[\dots] húng tôi cùng nhau mời cơm và bắt đầu bữa cơm trưa. Khi bữa ăn bắt đầu, mọi người
đều tò mò với bát nhỏ đặt trước mặt mỗi người. Trong đó chứa chất lỏng đen sánh, trông
có vẻ không giống bất kỳ món ăn truyền thống nào.

Noel-chan nhíu mày nhìn bát mật, quay sang hỏi tôi: ``Đây là cái gì vậy?'' Tôi đáp gọn:
``Mật mía đấy.''

Roboco-chan nghe vậy, ngơ ngác: ``Mật mã? Nhưng nó có ăn được đâu?'' Tôi bật cười: ``Anh nói `mật mía', không phải `mật mã'. Có cần anh khám lại thính giác cho em không đây, Ngố ơi?''

Nghe xong, nàng Người máy phồng má lên, đỏ mặt: ``Đừng có mà trêu vậy trước mặt mọi
người chứ!''

Haato-chan xen vào, nghiêng đầu nhìn chất lỏng sánh trong bát của mọi người: ``Nhìn nó
giống mật ong nhỉ?'' Mẹ tôi ngồi đối diện chúng tôi, giải thích: ``Gần giống như vậy, nhưng mà mật mía sánh
hơn.''

Subaru-chan gật gù: ``Chắc nó ngọt hơn mật ong nữa, đúng không?'' Tôi đưa mắt ra hiệu:
``Cứ thử đi, mỗi người một bát nhỏ. Anh dám chắc đây sẽ là một trải nghiệm khó quên
đấy.''

Ban đầu, tất cả mọi người nhìn chiếc bát nhỏ ấy với đầy rẫy sự nghi ngờ. Nhưng rồi,
Miko-chi can đảm (?), tò mò chấm đầu đũa vào bát mật và nếm thử. Vừa cho vào miệng,
mặt em ấy khẽ nhăn lại: ``Ngọt quá ne! Ngọt hơn cả kem bánh taiyaki nữa!''

Nhìn vẻ mặt không thoải mái của em ấy, tôi phì cười: ``Thế nên nó mới ăn kèm với bánh
chưng. Anh biết mấy đứa cũng bắt đầu ngán bánh chưng ăn mặn rồi, đúng không?''

Mel-chan và A-san thì bát mật, rồi nàng Dơi lên tiếng tỏ vẻ ngần ngại: ``Nhưng mà liệu
đồ mặn với đồ ngọt trộn với nhau có ổn không?''

Tôi mỉm cười, đáp: ``Nhiều thứ tưởng không liên quan với nhau, nhưng khi kết hợp lại
thành siêu phẩm đấy.'' Thậm chí, A-san cũng phải gật gù đồng tình: ``Ẩm thực nơi này có
nhiều điều bất ngờ lắm đấy. Nó cũng đáng thử.''

\il

Mel-chan ngồi bên cạnh tôi, ngạc nhiên thốt lên khi nhìn thấy Noel-chan ăn lấy ăn để,
như thể em ấy đang bị bỏ đói: ``Chị vẫn không thể tin là Noel-chan ăn nhiều đến thế.''

Nàng Hiệp sĩ ra vẻ tự hào với ``thành tích'' ăn gần như không ai có thể vượt qua được:
``Có thịt là em ăn thôi!''

A-san lên tiếng nhắc nhở: ``Coi chừng béo lên đấy, Noel. Đừng quên rằng cô là một idol
nữa đó.''

Pekora-chan chen vào, nhại lại giọng điệu của chị ấy: ``Đúng rồi, peko! Đừng để lúc nào
lên sân khấu cũng thấy một quả bóng lăn qua lăn lại nhé.''

Nàng Hiệp sĩ giả vờ nhăn nhó với những lời nhắc nhở từ hai người họ: ``Pekora-chan ác
quá! Mình không có béo đến thế đâu!''

Tôi bật cười khi nhớ lại lúc trong phim, cả nhóm còn thi nhau thử sức với món bánh
chưng. Matsuri-chan đã lỡ tay làm rơi một miếng xuống đất, còn Shion-chan thì khẳng
định rằng mình không ăn được hành dù chẳng ai ép.

Nàng Dơi cười khúc khích, nhớ lại những kỷ niệm này: ``Chị nhớ lúc đó Matsuri-chan
còn giả vờ là mình không làm rơi miếng bánh chưng ấy.''

Nghe vậy, người mà em ấy nhắc tên liền thốt ra lời phản đối: ``Thì có ai bắt bồi thường
đâu, tự dưng mọi người cứ nhìn chằm chằm tôi đấy thôi!''

Mio-chan thêm vào, cà khịa: ``Chắc là vị chị làm rơi nó trước mặt camera đấy chứ.''

Nàng Lễ hội đỏ mặt (lần thứ ba), chống ché: ``Thế thì đừng đưa cảnh đó vào phim chứ!
Kuon-senpai! A-chan! Mau cắt cảnh đó đi!''

Cả nhóm phá lên cười, còn tôi chỉ biết lắc đầu, cố gắng kìm lại tiếng cười của mình.

Cảnh phim chuyển sang buổi sáng mùng Ba, khi các cô gái \hll\ tập trung trong một lớp học tại trường tôi theo học để làm bài kiểm tra toán mà chính tay tôi đã chuẩn bị. Không khí trở nên căng
thẳng hơn hẳn so với những cảnh trước. Từng người một cầm bút, trán lấm tấm mồ hôi,
trong khi tôi đứng phía trước, đóng vai một giám thị nghiêm khắc.

\il

[\dots] au khi Kuon tuyên bố bắt đầu làm bài, không khí trong phòng thi trở nên yên ắng đến lạ
thường. Tiếng lật giấy loạt xoạt và tiếng bút chì kẽ trên phiếu trả lời trắc nghiệm là những âm thanh duy nhất vang lên. Dù vậy, sự căng thẳng hiện rõ trên từng khuôn mặt.

Fubuki và Miko, sau màn ngất xỉu bất đắc dĩ, được mọi người cố gắng lay gọi. Đến khi
tỉnh lại, cả hai vẫn trong trạng thái hoảng loạn:

\begin{dialogue}
  \speak{Miko} Tại sao đề toán này khó thế, ne!?
  \speak{Fubuki} Em thậm chí còn không nhớ lần cuối mình làm bài toán là khi nào nữa\dots
\end{dialogue}

Sora nắm lấy tay Miko, nói nhỏ, giọng an ủi: ``\hl{Không sao đâu, làm hết sức mình là đượ.
  Mình tin là cậu có thể làm được, Miko-chan!}''

Mel đứng bên cạnh cũng gật đầu, nhẹ nhàng động viên: ``Đúng đấy! Với lại, đây là bài
trắc nghiệm mà, cứ chọn bừa cũng được!''

Những lời an ủi của họ cuối cùng cũng giúp Miko và Fubuki trấn tĩnh hơn một chút, đủ
để bắt đầu cầm bút và thử làm bài. Dù vậy, cứ vài phút, tiếng than thở của nàng Vu nữ
lại vang lên: ``\textit{Không hiểu sao bài này lại phải tính cả cái tích phân lằng nhằng này chứ
  ne\dots\ Mấy cái này thì giúp gì được trong đời thực cơ chứ, ne!}''

Ở phía đầu lớp, Kuon, Choco và A-chan đã chuẩn bị sẵn máy quay để ghi lại khoảnh khắc
đặc biệt này. Nàng Đeo kính cầm máy quay, còn cậu Tổng đứng bên cạnh, vừa nhìn vừa
chỉ tay, không quên thêm lời nhận xét: ``Chị thấy nét mặt của Subaru-chan không? Nhìn
có khác gì đang gặp ác mộng đâu.''

A-chan khẽ bật cười, hướng máy quay về phía nàng Vịt, người đang vò đầu bứt tai khi
cố gắng giải một bài toán tổ hợp: ``Chị nghĩ phải quay lại cảnh này. Mấy đứa về xem lại
chắc sẽ tự cười bản thân đấy.''

Choco cũng bật cười khi nhìn vẻ mặt lúng túng của tất cả mọi người: ``Đề thi mà anh cho
đúng là dị thật đấy\dots\ Không biết họ sẽ xử trí thế nào nhỉ.''

Korone ngồi cách đó không xa, đôi tai chó của cô khẽ cụp xuống, vừa nhai bút vừa nhìn
chằm chằm vào câu hỏi: ``\textit{Hmm\dots\ Nếu chọn sai thì chọn đáp án nào nhỉ? A hay B nhỉ\dots?
  Hay cứ C luôn cho chắc?}''

Noel ở phía bên kia phòng thì đập tay lên bàn khe khẽ, miệng lẩm bẩm: ``\textit{Hình như mình
  nhớ công thức này\dots\ Không đúng, sai rồi! Công thức nào mới đúng nhỉ!?}''

Cả phòng thi lúc này, dù im lặng, nhưng ai cũng toát lên vẻ căng thẳng. Chỉ có Kuon,
đứng một góc, vừa liếc nhìn đồng hồ đeo tay của mình, vừa mỉm cười như thể rất hài lòng
với bài kiểm tra đầy thách thức của mình.

\il

Subaru-chan xem đến đoạn mình gục đầu xuống bàn trong phim thì ôm mặt, thở dài: ``Đề
của anh đúng là vừa khó\dots\ vừa ác liệt\dots''

Rushia-chan đỏ mặt, khẽ gật đầu đồng tình: ``Lại còn thao túng tâm lý bọn em nữa chứ!''

Tôi cười, giải thích thêm khi mọi người quay sang nhìn mình: ``Thế mới thấy không khí
kỳ thi của bọn anh căng thẳng đến mức nào! Cũng may các em là idol đấy, chứ nếu không
chắc chẳng ai chịu nổi đâu.''

Akutan xem đến cảnh mình bấm máy tính và gãi đầu, bật cười chua xót: ``Anh nói cũng
phải\dots\ Nhưng mà bài này đúng là khó thật!''

Fubuki thì chỉ vào đoạn phim vào bản thân đang tính nhẩm bằng tay do máy tính bị hỏng,
lắc đầu than: ``Toán vốn dĩ có phải là sở trường của em đâu mà\dots\ May mà em không bị
trượt\dots''

Tôi nhún vai trong thực tế, đáp lại với giọng nửa đùa nửa thật: ``Chỉ cần không ai gục ngã
giữa chừng là anh đã mừng lắm rồi.''

Ayame-chan phồng má lên, dường như vẫn còn hơi hậm hực từ vụ bài kiểm tra đó: ``Em
vẫn chưa bỏ qua việc anh bắt em nộp bài khi chỉ còn mười câu nữa thôi đấy!''

Tôi mỉm cười, chỉ về phía nàng Hiệp sĩ: ``Còn đỡ cay cú hơn ai đó thiếu đúng một câu là
thoát phạt đấy, Ojou-san ạ.''

Peko-chan xem đến cảnh mình đang reo hò bêu xấu hai người điểm thấp nhất trong bài
kiểm tra đó thì phá lên cười: ``Đoạn chúng ta diễu hành khắp trường đúng là vui thật đó,
peko!''

Cả Noel-chan và Roboco-chan đều cùng lắc đầu ngán ngẩm, dường như muốn quên đi
``vết nhơ'' đó: ``Vui cái nỗi gì chứ\dots''

A-san thở dài, ánh mắt vừa bất lực vừa buồn cười: ``Cũng may họ cũng là các idol. Chứ
không chắc họ cũng cơm toi luôn rồi.''

Tôi gãi đầu, cười trừ: ``Em nghĩ các em ấy cần học thêm chút logic đấy.''

Khi cảnh quay kết thúc, cả nhóm cùng bật cười vì những biểu cảm kỳ lạ của mình trên
màn hình, đặc biệt là chính bản thân hai ``phạm nhân'' mà các cô gái của chúng ta đi kêu
la khắp trường.

Cảnh phim chuyển tới đoạn lễ hội, nơi các cô gái tham gia vào không khí rộn ràng và
náo nhiệt. Những quầy hàng đầy màu sắc, tiếng trống giòn giã, và đặc biệt là màn múa
lân phun lửa của Ayame-chan. Bộ tứ FAMS đảm nhận vai chính trong điệu múa lân, với
Miko-chan chỉ huy, còn Okayu-chan và Korone-chan làm tay trống, tạo nên một màn biểu
diễn đầy sôi động và ấn tượng.

\il

[\dots] Tiếng trống càng lúc càng dồn dập, tạo nên nhịp điệu sôi động. Trong khi đó, Miko, đứng
ở vị trí chỉ huy phía trước, hăng say vung gậy trừ tà. Nhưng rõ ràng, sự tự tin của cô nàng
không đi kèm với kỹ thuật thuần thục. Cô nàng lẩm bẩm, tay vẫn cầm cây gậy không
ngừng chỉ tứ tung: ``\textit{Như thế này\dots\ đúng không nhỉ, ne?}''

Từ phía khán giả, Kuon nhếch miệng cười, nói nhỏ với A-chan: ``{Em ấy trông có vẻ làm
  tốt đấy\dots}''

Nàng Đeo kính khẽ gật đầu: ``Công nhận. Nhưng mà\dots\ trông hơi lạ.''

Không những vậy, trong lúc Miko hăng hái chỉ huy, cô còn vô tình tạo ra một số tư thế
đặc biệt. Những động tác này chẳng những thu hút ánh mắt ngạc nhiên của khán giả mà
còn khiến các đội múa lân khác phải há hốc mồm kinh ngạc.

Dường như không để ý tới đội múa của mình, Miko xoay gậy vòng quanh, tạo nên một
vòng cung lớn, rồi hét lên: ``\textbf{Tiến lên! Mọi người nhảy đi nào!}''

Như thể bị lây năng lượng của nàng Vu nữ, Ayame hô lớn: ``\textbf{A-lê-hấp!}'' Con lân đột nhiên
nhảy vọt lên theo hiệu lệnh của Miko, làm cả đám đông phải bật cười, đồng thời cổ vũ
một cách nhiệt liệt.

Subaru ở phía sau, không giấu được vẻ ngỡ ngàng trước hành động bất ngờ của cô bạn,
hét toáng lên: ``\textbf{Này- Bà làm cái gì vậy!?}''

Tiếng trống lại vang lên dồn dập hơn, như tiếp thêm lửa cho sự náo nhiệt của màn trình
diễn. Nhưng những động tác bất ngờ của nhóm FAMS khiến cả khán giả và những nhóm
đang trình diễn khác đều không thể rời mắt.

\il

Matsuri-chan vỗ tay hào hứng khi đoạn quay màn múa lân xuất hiện: ``Quả nhiên lễ hội
vẫn là vui nhất!''

Shion-chan nhìn cảnh nàng Oni làm cú nhảy xoay người trong bộ lân, rồi khoanh tay, tỏ
vẻ hờ hững: ``Đúng với chị thôi, chứ em thấy chán òm.''

Aqua-chan bỗng nhiên chỉ vào màn hình, nơi mà chính nàng Phù thủy đang mỉm cười giơ
một con tò he hình vịt: ``Nghe ai đó nặn ra con tò he hình vịt nói gì kìa\~{}''

Choco-sensei bật cười nhẹ, cố xoa dịu không khí: ``Thôi nào, đó cũng là trải nghiệm văn
hóa truyền thống chứ?''

Haato giơ lấy tấm giấy đỏ mà em ấy đã viết từ ngày hôm đó ra, reo lên: ``Chị vẫn còn giữ
chữ `Phúc' mà chị viết ngày hôm đó này!!''

``Giờ không phải là lúc khoe đâu, Haato-senpai\dots'' Ru-chan liếc cô nàng đang hăng máu
đó.

Fubuki-chan nghiêng đầu nhìn cảnh mình cười rạng rỡ khi được trao một chiếc lồng đèn
nhỏ, cười khì: ``Chẳng ai để ý là chị suýt làm rơi lồng đèn đâu nhỉ.''

Okayu vừa nhai khoai lang nướng (từ một xó nào đó) vừa nói, mắt vẫn dán vào màn hình:
``Còn phần trống của bọn em, nghe đỉnh thật đấy. Koro-san thì hơi lệch nhịp chút thôi.''

Nàng Khuyển giật mình quay sang cô bạn, tỏ vẻ bất mãn tới độ đập nhẹ vào vai: ``Này
Ogayu, đừng có đổ lỗi hết cho tôi nhé! Kuon-senpai bảo là cứ đánh thoải mái rồi mà!''

Tôi cười trừ khi bị nhắc đến, khẽ gật đầu: ``Ừ thì\dots\ lễ hội mà, tự do sáng tạo chút cũng
không sao. Miễn là mọi người thấy hay là được mà, nhỉ?''

A-san nhìn tôi, lắc đầu ngán ngẩm nhưng lại mỉm cười: ``Hikawa-san đúng là luôn tìm đủ
lý do để bào chữa cho mọi người nhỉ.''

Subaru, trong cảnh quay, bật cười khi thấy mình lúng túng cầm chiếc đầu lân quá nặng, đỏ mặt phân bua: ``Nặng thật đấy! Ai bảo không có phần hướng dẫn chi tiết nào cho chúng
tôi cả!''

Miko-chi xem lại màn trình diễn của mình với mọi người, tỏ vẻ tự hào về vai trò của mình:
``Nhưng mà nhờ chị, màn múa lân này mới thành công đấy ne!''

``Và cũng chính người đó đã rút ra đúng quẻ `Đại Hung' đấy\~{}'' Aki châm chọc nói.

``I-Im đi.''

Ayame-chan nhướng mày khi nghe vậy, bật cười khúc khích: ``Màn phun lửa cuối cùng do
mình làm ấy\dots\ Chắc mọi người cũng nên cảm ơn mình chứ nhỉ, yo?''

``Đến mức còn phải chườm đá là đủ hiểu rồi\dots'' nàng Phù thủy bông đùa. Cả nhóm bật
cười khi đoạn phim kết thúc với cảnh mọi người cùng hò reo với màn trình diễn siêu đỉnh
như vậy.

Bộ phim dần đi đến hồi kết khi khung cảnh lễ hội nhường chỗ cho một cảnh quay toàn bộ
các cô gái cùng gia đình tôi đứng trên tầng thượng, ngắm nhìn ánh sáng rực rỡ từ pháo
hoa đêm 30 Tết. Âm nhạc lắng dần, chuyển sang một giai điệu nhẹ nhàng, ấm áp.

Hình ảnh chuyển sang những cảnh đằng sau hậu trường: Subaru-chan đang cố giơ đầu lân
nhưng lại bị mất thăng bằng ngã sõng soài; Ayame-chan thì cười phá lên khi thấy cặp đôi
Chó-Mèo tranh cãi vì ai đã đánh trống lệch nhịp; Shion-chan cố gắng thổi bóng nhưng
làm bóng bay nổ ngay trên mặt, khiến mọi người cười lăn.

Chuyển tới một cảnh khác, Fubuki-chan bị một con mèo vô tình nhảy lên người khi đang
quay, làm cô hét lên rồi cười ngượng; Noel-chan và Choco-san cố gắng giấu đi sự bối rối
khi làm rơi đĩa bánh ra sàn; còn tôi, trong một cảnh quay thử, đứng điều chỉnh máy quay
nhưng lại bị chính nó\dots\ ngã đập vào trán.

Sau khi những tiếng cười vang lên từ hậu trường, màn hình trở lại với cảnh toàn đoàn
cùng nhau cúi chào dưới bầu trời đầy sao. Một lời bình đơn giản vang lên:

\begin{center}
  \uline{Và đó là cách mà các thành viên \hll\ trải nghiệm Tết ở Etsunan. Không biết trong
    tương lai họ sẽ trải nghiệm thêm những gì nhỉ?}
\end{center}

Sau đó, đoạn chạy chữ (credits) hiện lên, với danh sách toàn bộ các thành viên \hll\
A-san, cùng tên của gia đình hai bên nội ngoại nhà tôi. Những hình ảnh nhỏ xuất hiện ở
góc màn hình, ghi lại thêm các khoảnh khắc vui nhộn: Choco-chan dọn bánh rơi, Ru-chan
đỏ mặt vì bị chọc ghẹo, hay Pekora-chan cố giành lì xì của người khác nhưng bị bắt quả
tang.

Âm nhạc kết thúc bằng một giai điệu tươi sáng, khép lại bộ phim đầu tiên của \hll\ tại
Etsunan.

Cuối cùng, một chữ ``おわり (HẾT)'' rất lớn hiện ra ở giữa màn hình, đi kèm với đó là
những cánh hoa dào đang rơi lả tả xuống làm nền, báo hiệu rằng một mùa Tết đã qua đi.

\ct

Sau khi bộ phim kết thúc, căn phòng rơi vào một khoảng lặng ngắn. Tất cả mọi người
dường như vẫn đang ngẫm nghĩ về những gì vừa xem, cho đến khi Mio-chan lên tiếng,
phá vỡ bầu không khí: ``Em thấy nó hơi bất ổn\dots'' em ấy tỏ vẻ nghiêm túc, ánh mắt không
rời màn hình vừa tắt.

``Tất nhiên rồi,'' tôi đáp, nhún vai, ``đây chỉ là những cảnh chính mà anh định cho mấy
đứa xem trước thôi mà.''

Sói đen-chan nhíu mày: ``Nhưng em thấy có nhiều cảnh bị cắt quá! Anh không sợ như
thế sẽ khiến người ta hiểu nhầm sao?''

Mel-chan ngồi bên cạnh, gật đầu hưởng ứng: ``Đúng đó! Cả bộ chỉ dài có tám phút thôi.
Em tưởng các anh định làm một bộ dài 90 phút cơ chứ!''

Fubu-chan cũng lên tiếng: ``Thậm chí lên tiếng hai hay ba tiếng lận chứ! Thế này chẳng
bõ chút nào!''

Tôi bật cười: ``Bọn anh làm theo từng tập ngắn thôi. Mỗi tập như thế sẽ phù hợp hơn với
thời lượng mà mọi người có thể xem trực tuyến.''

Dơi-chan chớp mắt đầy thắc mắc: ``Ể? Anh nói tập để xem trực tuyến là sao?''

``Chẳng phải bây giờ đại dịch COVID đang hoành hành à?'' Tôi giải thích, giọng nghiêm
túc hơn. ``Rạp chiếu phim đóng cửa hết cả rồi, làm phim dài thì cũng không chiếu được.
Mọi người chỉ có thể xem qua mạng thôi.''

A-san gật đầu đồng tình: ``Đúng thế. Giờ mà làm phim chiếu rạp thì có ai đến xem đâu.
Thế nên bọn tôi mới phải chuyển sang định dạng như thế này.''

Aqua-chan ngồi phía cuối hàng ghế, ngạc nhiên thốt lên: ``Ồ\dots\ giờ thì em hiểu rồi.'' Cáo
trắng-chan vẫn giữ vẻ điềm tĩnh, khẽ cười: ``Chẳng trách mà em thấy ngoài đường vắng
tanh đến thế. Ai cũng ở nhà hết cả.''

Rushia nghe em ấy nói như vậy, lo lắng hỏi: ``Thế ở thành phố của anh có bị ảnh hưởng
nhiều không?''

``Có chứ!'' tôi thở dài, nhìn quanh cả nhóm. ``Anh còn phải chuẩn bị tinh thần học và làm
việc tại nhà đây này. Thế nên anh cũng khuyên mấy đứa cẩn thận, đừng ra ngoài nhiều
quá.''

Không khí trở nên chùng xuống trong giây lát. Nhưng rồi tôi vỗ tay, kéo mọi người trở lại
chủ đề chính: ``Nói thế thôi, giờ quay lại vấn đề quan trọng nào. Mấy đứa thấy bộ phim
thế nào?''

Mọi người đứng lại một lúc lâu, có vẻ ai cũng đang suy nghĩ về bộ phim vừa xem. Miochan
tiếp tục là người phá vỡ sự im lặng đầu tiên, vẫn giữ dáng vẻ nghiêm túc: ``Em thấy
mọi thứ rất lộn xộn.''

Pekora-chan bật cười khúc khích, nhanh chóng đáp lại: ``Giống như lúc chúng ta ở văn
phòng ấy.''

Rushia gật đầu đồng tình, mặt hơi cau có: ``Chẳng khác gì những trò ngớ ngẩn chúng ta
hay làm hàng ngày cả.''

Tôi mỉm cười. Thấy rằng mọi người đang dần có phản ứng cùng với ba cô gái kia, tôi đáp
lại với giọng trầm ấm nhưng không kém phần tự tin: ``Đó mới là điều anh hướng tới: cho
mọi người biết \hll\ thực sự như thế nào.''

``Ủa?'' tất cả mọi người ngạc nhiên, đôi mắt mở lớn như muốn hiểu rõ hơn ý của tôi. ``Ý
anh là sao chứ, Kuon-senpai?''

Tôi đang định tiếp tục giải thích, thì đồng nghiệp của tôi xen vào, tiếp lời với giọng bình
thản nhưng rất rõ ràng: ``Bản chất \hll\ chính là những thứ gì đó tạp nham đến nhảm
nhí. Ngay cả \textit{hologra} cũng được xây dựng từ điều đó mà.''

\dots Đến cả tôi cũng bất ngờ với lời giải thích từ một người luôn phát điên với mọi trò nghịch
ngợm từ họ đấy.

Mel-chan nghiêng đầu, mắt sáng lên như vừa hiểu ra: ``Vậy nên anh chị mới làm một bộ
phim có nét tương đồng với cuộc sống \hll\ bình thường sao?''

``Đúng là như vậy,'' tôi gật đầu. ``Anh cũng muốn mọi người đổi gió một tí, để thoát khỏi
những thứ lặp đi lặp lại hàng ngày.''

Ayame-chan bật cười, tỏ vẻ thuyết phục với lời của chúng tôi: ``Đúng đó, cứ ở Nhật mãi
cũng chán. Thế nên là tôi đồng ý liền, yo!''

Aqua-chan lườm em ấy, vẻ không hài lòng: ``Bà thì dễ dãi quá rồi\dots'' Tôi khẽ bật cười
với phản ứng của nàng Hầu gái: ``Em cũng có vẻ rất vui mà. Anh nghĩ mọi người ai cũng
cảm thấy như vậy!''

A-san cười nhẹ, giọng chị tràn đầy suy tư: ``Tôi cũng muốn thấy mối quan hệ của mọi
người với nhau trong một môi trường mới sẽ ra sao. Có khi nhờ thế mà mọi người sẽ hiểu
nhau hơn.''

Không khí trở nên nhẹ nhàng và thân thiện hơn hẳn. Tất cả đều im lặng trong giây lát,
như để ngẫm nghĩ về điều mà chị ấy vừa nói.

Cuộc trò chuyện dần trở nên sôi nổi hơn sau những lời chia sẻ ban đầu. Fubuki-chan mỉm
cười, giọng đầy phấn khích: ``Nếu mà chị nói như vậy thì\dots\ Em nghĩ mọi người làm khá
tốt đấy!''

Aki-chan dịu dàng tiếp lời: ``Chưa bao giờ mình thấy mọi người vui vẻ với nhau như vậy
á.''

Noel-chan thì thản nhiên góp ý, ánh mắt lấp lánh: ``Cũng chưa bao giờ em được ăn nhiều
của ngon vật lạ đến thế!''

``Tham ăn vừa thôi, peko,'' Peko-chan lập tức nhắc nhở, khiến mọi người bật cười.

Tôi nhìn quanh nhóm một lượt rồi hỏi: ``Nếu mấy đứa thấy vui thì chắc chắn người xem
cũng sẽ thấy vui thôi. Anh hỏi thật nhé, các em có thấy thoải mái khi ăn Tết ở quê anh
không?''

Shion-chan hơi ngập ngừng, nhưng cuối cùng cũng thốt ra một câu: ``\dots Một chút.'' Akutan
cười tươi rói: ``Thấy chưa, tôi đã bảo rồi mà\~{}''

Ngược lại, Roboco-chan cười tươi rói: ``Chị thấy vui mà. Chưa bao giờ chị được chơi đã
như vậy!''

Ayame-chan giơ hai tay hưởng ứng: ``Trừ bài kiểm tra ra, còn lại thì cái nào cũng vui cả,
yo! Em thích nhất là lúc mà bọ em đi múa lân!''

Sora-chan nghiêng đầu, ánh mắt như đang mơ màng nhớ lại: ``Nó còn khác hẳn so với
những gì em nghĩ nữa.''

Mat-chan thì thẳng thắn hơn: ``Em muốn được trải nghiệm Tết một lần nữa!''

Tôi khẽ gật đầu, cảm thấy nhẹ nhõm: ``Vậy là anh vui rồi. Mấy đứa muốn làm gì tiếp theo
đây?''

Tưởng chừng câu hỏi đó sẽ mở ra một chủ đề khác, nhưng không ngờ tất cả mọi người
lại đồng thanh thốt lên, một số người trong số họ thậm chí còn sáng mắt lên: ``\textbf{Tất nhiên
  là được anh lì xì lại cho chúng em rồi!}''

Câu trả lời đầy nhiệt tình (một số còn thái quá) làm tôi khựng lại, sững sờ trong giây lát
trước khi đổ mồ hôi hột: ``Anh chưa chuẩn bị tiền đâu\dots''

Tiếng cười giòn tan vang lên, lấp đầy không khí ấm áp của buổi gặp mặt. Ai nấy đều tươi
cười trước phản ứng bất ngờ của tôi, và khung cảnh ấy như một dấu ấn đáng nhớ cho cái
Tết đặc biệt của chúng tôi bên nhau.

``Vậy mấy đứa có lời gì muốn chúc nhau trong năm Canh Tý sắp tới không?'' tôi hỏi.
``Anh biết rằng một năm sắp tới sẽ là một năm đầy khó khăn gian khổ\dots''

``\dots nhưng chúng ta vẫn phải cố gắng thôi!'' A-san tiếp lời. ``Tôi muốn thấy hoài bão của
mọi người trong năm mới này!''

Sora-chan đứng lên, nở nụ cười rạng rỡ như ánh nắng đầu xuân, vẫy tay với mọi người.
``Mọi người à, ta cùng cố gắng phấn đấu cho một năm mới sắp tới nha!''

Tôi không thể nhịn được mà bật cười, nhắc nhẹ: ``\dots Dù nó đã qua một tháng rồi.''

Em ấy không hề tỏ vẻ sượng trân với lời nhắc đó, dứt khoát đáp lại: ``Cố lên nào, mọi
người!''

Cảm hứng từ em lan tỏa nhanh chóng khắp phòng, tất cả mọi người đồng thanh giơ tay
hô lớn: ``Cố lên!''

Một năm mới âm lịch đã bắt đầu, mang theo hy vọng, thách thức và cả những chông gai
không thể lường trước. Trong bối cảnh đại dịch toàn cầu vẫn còn, mỗi người đều gánh
trên vai trách nhiệm, không chỉ với bản thân mà còn với cả cộng đồng.

\hll, với sức mạnh từ tình bạn, sự đoàn kết, và những tiếng cười, sẽ đối mặt với
những thử thách đó ra sao?

Chỉ thời gian và nghị lực của họ mới có thể trả lời được.

Với tất cả mọi người, tôi xin gửi lời chúc:

\begin{center}
  \uline{\textbf{Chúc mừng năm mới năm Canh Tý!}}
\end{center}

\end{document}